% ==============================================================================
% MÉMOIRE DE FIN D'ÉTUDES — SYSTÈME D'ANALYSE DE SENTIMENTS
% Auteur : Hissein Nasser
% Compilateur : pdflatex (2 passes minimum)
% ==============================================================================
\documentclass[12pt,a4paper,oneside]{report}

% ── Encodage et langue ─────────────────────────────────────────────────────────
\usepackage[utf8]{inputenc}
\usepackage[T1]{fontenc}
\usepackage[french]{babel}
\usepackage{lmodern}

% ── Mise en page ───────────────────────────────────────────────────────────────
\usepackage[top=2.5cm, bottom=2.5cm, left=3cm, right=2.5cm]{geometry}
\usepackage{setspace}
\onehalfspacing

% ── Couleurs ───────────────────────────────────────────────────────────────────
\usepackage{xcolor}
\definecolor{bleuprincipal}{RGB}{31,73,125}
\definecolor{bleuclaire}{RGB}{46,134,193}
\definecolor{vertprincipal}{RGB}{31,138,112}
\definecolor{rougeprincipal}{RGB}{231,76,60}
\definecolor{griscode}{RGB}{248,249,250}
\definecolor{grisbordure}{RGB}{204,204,204}
\definecolor{orangeaccent}{RGB}{211,84,0}
\definecolor{jauneneutre}{RGB}{241,196,15}

% ── Graphiques et figures ──────────────────────────────────────────────────────
\usepackage{graphicx}
\usepackage{float}
\usepackage{caption}
\usepackage{subcaption}
\captionsetup{font=small, labelfont={bf,color=bleuprincipal}}

% ── Tableaux ───────────────────────────────────────────────────────────────────
\usepackage{booktabs}
\usepackage{tabularx}
\usepackage{multirow}
\usepackage{array}
\usepackage{longtable}
\usepackage{colortbl}

% ── TikZ et diagrammes ────────────────────────────────────────────────────────
\usepackage{tikz}
\usepackage{pgfplots}
\pgfplotsset{compat=1.18}
\usetikzlibrary{shapes.geometric, arrows.meta, positioning, calc,
                shadows, decorations.pathreplacing, fit, backgrounds,
                matrix, patterns, mindmap}

% ── Listings (code source) ────────────────────────────────────────────────────
\usepackage{listings}
\usepackage{lstautogobble}

\lstdefinestyle{stylePython}{
  language=Python,
  backgroundcolor=\color{griscode},
  basicstyle=\ttfamily\footnotesize,
  keywordstyle=\color{bleuprincipal}\bfseries,
  stringstyle=\color{rougeprincipal},
  commentstyle=\color{vertprincipal}\itshape,
  numberstyle=\tiny\color{gray},
  numbers=left,
  numbersep=8pt,
  breaklines=true,
  breakatwhitespace=false,
  frame=single,
  rulecolor=\color{grisbordure},
  tabsize=4,
  showstringspaces=false,
  captionpos=b,
  autogobble=true,
  xleftmargin=1cm,
  morekeywords={self, True, False, None, async, await, yield}
}

\lstdefinestyle{styleDart}{
  language=Java,
  backgroundcolor=\color{griscode},
  basicstyle=\ttfamily\footnotesize,
  keywordstyle=\color{bleuprincipal}\bfseries,
  stringstyle=\color{rougeprincipal},
  commentstyle=\color{vertprincipal}\itshape,
  numberstyle=\tiny\color{gray},
  numbers=left,
  numbersep=8pt,
  breaklines=true,
  frame=single,
  rulecolor=\color{grisbordure},
  tabsize=2,
  showstringspaces=false,
  captionpos=b,
  autogobble=true,
  xleftmargin=1cm,
  morekeywords={final, const, var, String, int, double, bool, void,
                Widget, StatelessWidget, StatefulWidget, BuildContext,
                Future, async, await, override, required, class}
}

\lstdefinestyle{styleJSON}{
  language=Python,
  backgroundcolor=\color{griscode},
  basicstyle=\ttfamily\footnotesize,
  keywordstyle=\color{bleuprincipal},
  stringstyle=\color{vertprincipal},
  numberstyle=\tiny\color{gray},
  numbers=left,
  numbersep=8pt,
  breaklines=true,
  frame=single,
  rulecolor=\color{grisbordure},
  captionpos=b,
  xleftmargin=1cm
}

\lstdefinestyle{styleSQL}{
  language=SQL,
  backgroundcolor=\color{griscode},
  basicstyle=\ttfamily\footnotesize,
  keywordstyle=\color{bleuprincipal}\bfseries,
  stringstyle=\color{rougeprincipal},
  commentstyle=\color{gray}\itshape,
  numberstyle=\tiny\color{gray},
  numbers=left,
  numbersep=8pt,
  breaklines=true,
  frame=single,
  rulecolor=\color{grisbordure},
  captionpos=b,
  xleftmargin=1cm
}

\lstset{style=stylePython}

% ── En-têtes et pieds de page ─────────────────────────────────────────────────
\usepackage{fancyhdr}
\pagestyle{fancy}
\fancyhf{}
\fancyhead[L]{\small\textcolor{bleuprincipal}{\leftmark}}
\fancyhead[R]{\small\textcolor{gray}{Analyse de Sentiments Étudiants}}
\fancyfoot[C]{\textcolor{gray}{\thepage}}
\fancyfoot[R]{\small\textcolor{gray}{Hissein Nasser — PFE 2026}}
\renewcommand{\headrulewidth}{0.5pt}
\renewcommand{\headrule}{\hbox to\headwidth{\color{bleuprincipal}\leaders\hrule height \headrulewidth\hfill}}
\renewcommand{\footrulewidth}{0pt}

% ── Style des chapitres ────────────────────────────────────────────────────────
\usepackage{titlesec}
\titleformat{\chapter}[block]
  {\normalfont\huge\bfseries\color{bleuprincipal}}
  {\colorbox{bleuprincipal}{\textcolor{white}{\quad\thechapter\quad}}}
  {1em}{}[\vspace{-0.5em}\rule{\textwidth}{2pt}]

\titleformat{\section}
  {\normalfont\Large\bfseries\color{bleuprincipal}}
  {\thesection}
  {0.8em}{}

\titleformat{\subsection}
  {\normalfont\large\bfseries\color{bleuclaire}}
  {\thesubsection}
  {0.8em}{}

\titleformat{\subsubsection}
  {\normalfont\normalsize\bfseries\color{orangeaccent}}
  {\thesubsubsection}
  {0.8em}{}

% ── Hyperliens ─────────────────────────────────────────────────────────────────
\usepackage[
  colorlinks=true,
  linkcolor=bleuprincipal,
  citecolor=vertprincipal,
  urlcolor=bleuclaire,
  pdfauthor={Hissein Nasser},
  pdftitle={Système d'Analyse de Sentiments des Commentaires Étudiants},
  pdfsubject={Mémoire de Fin d'Études},
  pdfkeywords={NLP, CamemBERT, Sentiment Analysis, FastAPI, Flutter, PyTorch}
]{hyperref}

% ── Bibliographie ──────────────────────────────────────────────────────────────
\usepackage[backend=bibtex,style=ieee,sorting=none]{biblatex}
\addbibresource{bibliography.bib}

% ── Acronymes ──────────────────────────────────────────────────────────────────
\usepackage[acronym,toc]{glossaries}
\makeglossaries

% ── Utilitaires ───────────────────────────────────────────────────────────────
\usepackage{amsmath}
\usepackage{amssymb}
\usepackage{enumitem}
\usepackage{mdframed}
\usepackage{pifont}
\usepackage{fontawesome5}
\usepackage{tcolorbox}
\tcbuselibrary{skins,breakable}

% ── Boîtes personnalisées ─────────────────────────────────────────────────────
\newtcolorbox{infobox}[1][]{
  enhanced, breakable,
  colback=bleuclaire!8!white,
  colframe=bleuclaire,
  fonttitle=\bfseries,
  title={#1},
  left=4pt, right=4pt, top=4pt, bottom=4pt
}

\newtcolorbox{definitionbox}[1][]{
  enhanced, breakable,
  colback=vertprincipal!8!white,
  colframe=vertprincipal,
  fonttitle=\bfseries,
  title={#1},
  left=4pt, right=4pt, top=4pt, bottom=4pt
}

\newtcolorbox{warningbox}[1][]{
  enhanced,
  colback=rougeprincipal!8!white,
  colframe=rougeprincipal,
  fonttitle=\bfseries,
  title={#1},
  left=4pt, right=4pt, top=4pt, bottom=4pt
}

% ── Commandes personnalisées ───────────────────────────────────────────────────
\newcommand{\code}[1]{\texttt{\textcolor{bleuprincipal}{#1}}}
\newcommand{\important}[1]{\textbf{\textcolor{rougeprincipal}{#1}}}

% ── Acronymes définis ─────────────────────────────────────────────────────────
% ==============================================================================
% ACRONYMES
% ==============================================================================
\newacronym{nlp}{NLP}{Natural Language Processing}
\newacronym{taln}{TALN}{Traitement Automatique du Langage Naturel}
\newacronym{bert}{BERT}{Bidirectional Encoder Representations from Transformers}
\newacronym{api}{API}{Application Programming Interface}
\newacronym{rest}{REST}{Representational State Transfer}
\newacronym{orm}{ORM}{Object-Relational Mapping}
\newacronym{sgbd}{SGBD}{Système de Gestion de Bases de Données}
\newacronym{ml}{ML}{Machine Learning}
\newacronym{dl}{DL}{Deep Learning}
\newacronym{ia}{IA}{Intelligence Artificielle}
\newacronym{http}{HTTP}{HyperText Transfer Protocol}
\newacronym{json}{JSON}{JavaScript Object Notation}
\newacronym{sql}{SQL}{Structured Query Language}
\newacronym{cors}{CORS}{Cross-Origin Resource Sharing}
\newacronym{ui}{UI}{User Interface}
\newacronym{ux}{UX}{User Experience}
\newacronym{gpu}{GPU}{Graphics Processing Unit}
\newacronym{cpu}{CPU}{Central Processing Unit}
\newacronym{ram}{RAM}{Random Access Memory}
\newacronym{pfe}{PFE}{Projet de Fin d'Études}
\newacronym{uml}{UML}{Unified Modeling Language}
\newacronym{mvc}{MVC}{Model-View-Controller}
\newacronym{rnn}{RNN}{Recurrent Neural Network}
\newacronym{lstm}{LSTM}{Long Short-Term Memory}
\newacronym{cnn}{CNN}{Convolutional Neural Network}
\newacronym{bpe}{BPE}{Byte-Pair Encoding}
\newacronym{mcd}{MCD}{Modèle Conceptuel de Données}
\newacronym{mld}{MLD}{Modèle Logique de Données}
\newacronym{crud}{CRUD}{Create, Read, Update, Delete}
\newacronym{ci}{CI}{Continuous Integration}
\newacronym{sdk}{SDK}{Software Development Kit}
\newacronym{ide}{IDE}{Integrated Development Environment}


% ==============================================================================
% DÉBUT DU DOCUMENT
% ==============================================================================
\begin{document}

\pagenumbering{gobble}

% ── Page de garde ─────────────────────────────────────────────────────────────
% ==============================================================================
% PAGE DE GARDE
% ==============================================================================
\begin{titlepage}
\thispagestyle{empty}

\begin{tikzpicture}[remember picture, overlay]
  % Bande supérieure bleue
  \fill[bleuprincipal] (current page.north west) rectangle
    ([yshift=-3.5cm]current page.north east);
  % Bande inférieure
  \fill[bleuprincipal!80!black] (current page.south west) rectangle
    ([yshift=1.8cm]current page.south east);
  % Bande latérale gauche
  \fill[vertprincipal] ([xshift=0cm]current page.north west) rectangle
    ([xshift=0.8cm]current page.south west);
  % Filet de séparation
  \fill[jauneneutre] ([xshift=0.8cm]current page.north west) rectangle
    ([xshift=1.1cm]current page.south west);
\end{tikzpicture}

\vspace*{0.5cm}
\begin{center}

% ── Logo / Établissement ──────────────────────────────────────────────────────
{\color{white}\Large\bfseries UNIVERSITÉ — DÉPARTEMENT INFORMATIQUE}\\[0.3cm]
{\color{white!80!gray}\large Filière : Génie Informatique}\\[0.2cm]
{\color{white!80!gray}\normalsize Option : Intelligence Artificielle \& Génie Logiciel}

\vspace{1.5cm}
\hrule height 1pt \relax
\vspace{0.3cm}
{\color{bleuprincipal!30!black}\small MÉMOIRE DE FIN D'ÉTUDES}\\
{\color{bleuprincipal!30!black}\small Pour l'obtention du diplôme de Licence / Master en Informatique}
\vspace{0.3cm}
\hrule height 1pt \relax

\vspace{1.8cm}

% ── Titre ─────────────────────────────────────────────────────────────────────
\begin{tcolorbox}[
  enhanced,
  colback=bleuprincipal!10!white,
  colframe=bleuprincipal,
  boxrule=2pt,
  arc=8pt,
  left=12pt, right=12pt, top=10pt, bottom=10pt,
  width=14cm
]
\begin{center}
{\huge\bfseries\color{bleuprincipal}
Système d'Analyse de\\[0.3em]
Sentiments des\\[0.3em]
Commentaires Étudiants\\[0.3em]
en Temps Réel
}\\[0.8em]
{\large\color{gray}
Basé sur le modèle CamemBERT, FastAPI, PostgreSQL et Flutter
}
\end{center}
\end{tcolorbox}

\vspace{2cm}

% ── Auteur et encadreur ───────────────────────────────────────────────────────
\begin{minipage}[t]{0.45\textwidth}
\begin{flushleft}
\textbf{\color{bleuprincipal}Réalisé par :}\\[0.5em]
\large Hissein Nasser\\
\small Étudiant en Informatique\\
\small hisseinnasser2003@gmail.com
\end{flushleft}
\end{minipage}
\hfill
\begin{minipage}[t]{0.45\textwidth}
\begin{flushright}
\textbf{\color{bleuprincipal}Encadré par :}\\[0.5em]
\large M./Mme. [Nom Encadreur]\\
\small Grade / Titre\\
\small Département Informatique
\end{flushright}
\end{minipage}

\vspace{2cm}

% ── Technologies ──────────────────────────────────────────────────────────────
\begin{tikzpicture}
  \node[fill=bleuprincipal!10, rounded corners=5pt, inner sep=6pt] at (0,0)
    {\small\color{bleuprincipal}\textbf{CamemBERT}};
  \node[fill=vertprincipal!15, rounded corners=5pt, inner sep=6pt] at (3,0)
    {\small\color{vertprincipal}\textbf{FastAPI}};
  \node[fill=rougeprincipal!10, rounded corners=5pt, inner sep=6pt] at (6,0)
    {\small\color{rougeprincipal}\textbf{PostgreSQL}};
  \node[fill=bleuclaire!15, rounded corners=5pt, inner sep=6pt] at (9,0)
    {\small\color{bleuclaire}\textbf{Flutter}};
  \node[fill=orangeaccent!10, rounded corners=5pt, inner sep=6pt] at (1.5,-0.9)
    {\small\color{orangeaccent}\textbf{PyTorch}};
  \node[fill=bleuprincipal!10, rounded corners=5pt, inner sep=6pt] at (4.5,-0.9)
    {\small\color{bleuprincipal}\textbf{Streamlit}};
  \node[fill=vertprincipal!15, rounded corners=5pt, inner sep=6pt] at (7.5,-0.9)
    {\small\color{vertprincipal}\textbf{SQLAlchemy}};
\end{tikzpicture}

\vspace{1.5cm}

% ── Année académique ──────────────────────────────────────────────────────────
{\color{white}\large\bfseries Année Académique 2025 -- 2026}

\end{center}
\end{titlepage}
\cleardoublepage


% ── Pages liminaires ──────────────────────────────────────────────────────────
\pagenumbering{roman}
\setcounter{page}{1}

% ==============================================================================
% DÉDICACE
% ==============================================================================
\chapter*{Dédicace}
\addcontentsline{toc}{chapter}{Dédicace}
\thispagestyle{empty}

\vspace*{3cm}

\begin{flushright}
\begin{minipage}{0.7\textwidth}
\begin{center}
{\Large\color{bleuprincipal}\textit{\og À ma famille, \fg{}}}

\vspace{1.5cm}

\normalsize\itshape
À mes parents, pour leur amour inconditionnel, leurs sacrifices
infinis et leur confiance inébranlable en moi, qui ont été
ma force et ma lumière tout au long de ce parcours.

\vspace{1cm}

À mes frères et sœurs, compagnons de vie, dont le soutien
moral et l'encouragement constant ont rendu ce travail possible.

\vspace{1cm}

À tous mes amis et collègues qui ont partagé ce chemin
académique avec moi, dans les moments de doute comme
dans les moments de réussite.

\vspace{1cm}

À tous ceux qui croient que la technologie, au service
de l'éducation, peut transformer le monde.

\vspace{2cm}

{\large\bfseries\color{bleuprincipal}--- Hissein Nasser ---}
\end{center}
\end{minipage}
\end{flushright}

\cleardoublepage

\input{frontmatter/remerciements}
% ==============================================================================
% RÉSUMÉ (FRANÇAIS)
% ==============================================================================
\chapter*{Résumé}
\addcontentsline{toc}{chapter}{Résumé}

\begin{tcolorbox}[
  enhanced,
  colback=bleuprincipal!5!white,
  colframe=bleuprincipal,
  boxrule=1.5pt, arc=6pt,
  left=10pt, right=10pt, top=8pt, bottom=8pt
]

Ce mémoire présente la conception et le développement d'un système complet
d'analyse automatique de sentiments appliqué aux commentaires d'étudiants
sur leurs cours universitaires. Face à la difficulté de traiter manuellement
des milliers de retours étudiants, ce projet propose une solution intelligente,
automatisée et temps réel, exploitant les avancées récentes du Traitement
Automatique du Langage Naturel (TALN).

\vspace{0.5em}

Le système repose sur le modèle de langage \textbf{CamemBERT}
(\textit{camembert-base}), une variante française de BERT développée par
Inria et CNRS, fine-tuné pour la classification de sentiments en trois
catégories : \textbf{Négatif}, \textbf{Neutre} et \textbf{Positif}. Le modèle
a été entraîné sur un jeu de données de 225 commentaires étudiants en langue
française, équilibrés entre les trois classes, avec les hyperparamètres
optimaux suivants : taux d'apprentissage de $2 \times 10^{-5}$, taille de
lot de 16, 3 epochs d'entraînement et une longueur maximale de séquence de
128 tokens. L'optimiseur AdamW avec un scheduler linéaire à palier de
chauffe a été utilisé pour garantir la convergence. Le modèle atteint
une précision de \textbf{79,41\%} sur le jeu de test.

\vspace{0.5em}

L'architecture globale du système comprend quatre composants principaux :
(1) un \textbf{backend RESTful} développé avec FastAPI et SQLAlchemy
communiquant avec une base de données PostgreSQL pour la persistance des
analyses ; (2) un \textbf{tableau de bord analytique} Streamlit enrichi
de visualisations interactives Plotly pour le suivi en temps réel des
tendances ; (3) une \textbf{application mobile} Flutter offrant aux
étudiants une interface moderne et intuitive pour soumettre leurs commentaires
et consulter les résultats ; (4) un \textbf{pipeline de prétraitement}
NLP complet comprenant la normalisation Unicode, la suppression des caractères
spéciaux et la tokenisation CamemBERT.

\vspace{0.5em}

Les résultats obtenus démontrent la faisabilité et l'efficacité d'une telle
approche, même avec un jeu de données limité, ouvrant la voie à des
améliorations significatives avec davantage de données annotées.

\end{tcolorbox}

\vspace{0.8em}
\noindent\textbf{Mots-clés :} Analyse de sentiments, TALN, CamemBERT,
Apprentissage automatique, FastAPI, Flutter, PostgreSQL, Classification de
texte, Fine-tuning, Deep learning.

\cleardoublepage

% ==============================================================================
% ABSTRACT (ENGLISH)
% ==============================================================================
\chapter*{Abstract}
\addcontentsline{toc}{chapter}{Abstract}

\begin{tcolorbox}[
  enhanced,
  colback=vertprincipal!5!white,
  colframe=vertprincipal,
  boxrule=1.5pt, arc=6pt,
  left=10pt, right=10pt, top=8pt, bottom=8pt
]

This thesis presents the design and development of a complete automated
sentiment analysis system applied to student feedback on university courses.
Faced with the challenge of manually processing thousands of student reviews,
this project proposes an intelligent, automated, and real-time solution
leveraging recent advances in Natural Language Processing (NLP).

\vspace{0.5em}

The system is built upon the \textbf{CamemBERT} language model
(\textit{camembert-base}), a French variant of BERT developed by Inria and
CNRS, fine-tuned for three-class sentiment classification: \textbf{Negative},
\textbf{Neutral}, and \textbf{Positive}. The model was trained on a dataset
of 225 French student comments, balanced across the three classes, using the
following hyperparameters: learning rate of $2 \times 10^{-5}$, batch size
of 16, 3 training epochs, and maximum sequence length of 128 tokens. The
AdamW optimizer with a linear warmup scheduler was employed to ensure
stable convergence. The model achieves an accuracy of \textbf{79.41\%} on
the test set.

\vspace{0.5em}

The overall system architecture consists of four main components:
(1) a \textbf{RESTful backend} developed with FastAPI and SQLAlchemy
interfacing with a PostgreSQL database for result persistence;
(2) an \textbf{analytics dashboard} built with Streamlit and enriched
with interactive Plotly visualizations for real-time trend monitoring;
(3) a \textbf{mobile application} built with Flutter providing students
with a modern, intuitive interface to submit comments and view analysis
results; (4) a complete \textbf{NLP preprocessing pipeline} including
Unicode normalization, special character removal, and CamemBERT
tokenization.

\vspace{0.5em}

The obtained results demonstrate the feasibility and effectiveness of such
an approach, even with a limited dataset, paving the way for significant
improvements with additional annotated data.

\end{tcolorbox}

\vspace{0.8em}
\noindent\textbf{Keywords:} Sentiment analysis, NLP, CamemBERT, Machine
learning, FastAPI, Flutter, PostgreSQL, Text classification, Fine-tuning,
Deep learning, Transfer learning.

\cleardoublepage


% ── Tables ────────────────────────────────────────────────────────────────────
{
  \hypersetup{linkcolor=black}
  \tableofcontents
  \listoffigures
  \listoftables
}

\printglossary[type=\acronymtype, title={Liste des Acronymes}]

\cleardoublepage
\pagenumbering{arabic}
\setcounter{page}{1}

% ── Corps du rapport ──────────────────────────────────────────────────────────
% ==============================================================================
% INTRODUCTION GÉNÉRALE
% ==============================================================================
\chapter*{Introduction Générale}
\addcontentsline{toc}{chapter}{Introduction Générale}
\markboth{INTRODUCTION GÉNÉRALE}{}

\vspace{0.3cm}

\begin{infobox}[Contexte de l'étude]
L'essor du numérique dans l'enseignement supérieur génère chaque année
des millions de retours étudiants. Traiter ces données manuellement est
devenu impossible. L'intelligence artificielle offre une réponse concrète.
\end{infobox}

\vspace{0.5cm}

\section*{Contexte général}

L'enseignement supérieur moderne est caractérisé par un flux continu et
croissant d'interactions entre enseignants et étudiants. Parmi celles-ci,
les commentaires et retours d'expérience sur les cours constituent une
source d'information précieuse pour améliorer la qualité pédagogique des
enseignements. Cependant, la multiplication des effectifs étudiants et la
diversification des modalités d'enseignement — notamment en ligne — ont
rendu le traitement manuel de ces retours pratiquement impossible à grande
échelle.

Face à ce constat, l'analyse automatique de sentiments (\textit{Sentiment
Analysis}) émerge comme une solution technologique particulièrement
prometteuse. Discipline à l'intersection du \acrfull{taln}, de l'apprentissage
automatique et de la linguistique computationnelle, elle permet d'identifier
et de classifier automatiquement les opinions exprimées dans un texte, sans
intervention humaine.

\section*{Problématique}

La langue française présente des défis spécifiques pour l'analyse
automatique de sentiments : richesse morphologique, usage fréquent de la
négation, expressions idiomatiques, niveaux de registre variés et présence
d'accents. Ces particularités rendent inadéquates les approches développées
pour l'anglais, et nécessitent des modèles spécialement adaptés à la langue
française.

Comment peut-on concevoir un système robuste, précis et accessible permettant
d'analyser automatiquement et en temps réel les sentiments exprimés dans les
commentaires d'étudiants en français, et de mettre ces analyses à la
disposition des enseignants et des administrateurs de manière exploitable ?

\section*{Solution proposée}

Ce projet propose une réponse complète à cette problématique à travers
le développement d'un système multi-couches intégrant :

\begin{itemize}[leftmargin=2em]
  \item[\textcolor{bleuprincipal}{\ding{118}}]
    Un modèle d'\acrfull{ia} basé sur \textbf{CamemBERT}, architecture
    Transformer pré-entraînée sur des milliards de tokens français, adaptée
    par fine-tuning à la classification de sentiments étudiants en trois
    catégories.

  \item[\textcolor{vertprincipal}{\ding{118}}]
    Un \textbf{serveur API RESTful} développé avec FastAPI, assurant
    l'interface entre le modèle \acrshort{ia} et les applications
    clientes, avec persistance des données dans PostgreSQL.

  \item[\textcolor{orangeaccent}{\ding{118}}]
    Un \textbf{tableau de bord analytique} Streamlit permettant aux
    enseignants de suivre en temps réel les tendances de sentiment par
    cours et par matière.

  \item[\textcolor{rougeprincipal}{\ding{118}}]
    Une \textbf{application mobile Android} Flutter offrant aux étudiants
    une interface simple et moderne pour soumettre leurs commentaires depuis
    n'importe où.
\end{itemize}

\section*{Objectifs du projet}

L'objectif général de ce \acrfull{pfe} est de concevoir, développer et
valider un système complet d'analyse de sentiments en temps réel pour les
commentaires étudiants en langue française. De manière plus spécifique,
ce projet vise à :

\begin{enumerate}[leftmargin=2.5em]
  \item Constituer et annoter un jeu de données de commentaires étudiants
        en français, équilibré entre les classes de sentiment.
  \item Fine-tuner le modèle CamemBERT pour la classification ternaire de
        sentiments avec des performances acceptables.
  \item Développer une API REST robuste et documentée exposant les
        fonctionnalités du modèle.
  \item Concevoir un tableau de bord analytique interactif pour la
        visualisation des résultats.
  \item Développer une application mobile multiplateforme permettant la
        soumission et la consultation des analyses.
  \item Évaluer les performances du système et identifier les pistes
        d'amélioration.
\end{enumerate}

\section*{Organisation du mémoire}

Ce mémoire est organisé en cinq chapitres, précédés d'une introduction et
suivis d'une conclusion générale :

\vspace{0.3em}
\noindent\textbf{Chapitre 1 — Contexte et Problématique :} présente le
contexte général de l'enseignement supérieur numérique, analyse les limites
des approches existantes et justifie les choix technologiques retenus.

\vspace{0.3em}
\noindent\textbf{Chapitre 2 — Analyse et Conception :} détaille l'analyse
fonctionnelle et non fonctionnelle du système, présente les diagrammes
\acrshort{uml} et l'architecture globale, et justifie les choix techniques.

\vspace{0.3em}
\noindent\textbf{Chapitre 3 — Implémentation :} décrit en détail les
environnements de développement, la structure du projet, les algorithmes
clés et les extraits de code commentés des composants principaux.

\vspace{0.3em}
\noindent\textbf{Chapitre 4 — Tests et Résultats :} présente la stratégie
de validation, les métriques de performance du modèle, les scénarios de
test fonctionnels et l'analyse critique des résultats obtenus.

\vspace{0.3em}
\noindent\textbf{Chapitre 5 — Discussion et Perspectives :} propose une
analyse critique du projet, identifie ses limites actuelles et trace les
perspectives d'évolution futures.

\cleardoublepage

% ==============================================================================
% CHAPITRE 1 : CONTEXTE ET PROBLÉMATIQUE
% ==============================================================================
\chapter{Contexte et Problématique}

\begin{infobox}[Objectif du chapitre]
Ce chapitre établit le cadre général de l'étude : le contexte éducatif
numérique, la problématique du traitement des retours étudiants, l'étude
des solutions existantes et la valeur ajoutée de notre approche.
\end{infobox}

% ──────────────────────────────────────────────────────────────────────────────
\section{Contexte général : La Révolution Numérique dans l'Éducation}
% ──────────────────────────────────────────────────────────────────────────────

L'enseignement supérieur traverse une transformation profonde depuis
l'avènement des technologies numériques. La massification de l'éducation
universitaire, conjuguée à l'explosion des plateformes d'apprentissage en
ligne (LMS — \textit{Learning Management Systems}), a multiplié
exponentiellement les interactions numériques entre enseignants et étudiants.

Selon plusieurs études pédagogiques récentes, un étudiant génère en moyenne
plusieurs dizaines de commentaires, évaluations et retours textuels par
semestre académique. À l'échelle d'un département ou d'une université,
cela représente des milliers, voire des dizaines de milliers de textes à
analyser chaque année.

\subsection{L'importance du retour étudiant}

Les retours des étudiants constituent un indicateur fondamental de la
qualité pédagogique. Ils permettent aux enseignants d'identifier les points
faibles de leurs cours, d'adapter leurs méthodes pédagogiques et d'améliorer
l'expérience d'apprentissage. Pour les administrations universitaires, ces
données agrégées fournissent une vision macroscopique de la satisfaction des
apprenants et orientent les décisions d'investissement pédagogique.

Cependant, la valeur réelle de ces données ne peut être exploitée que si
elles sont analysées de manière systématique, objective et rapide. C'est
précisément là que réside le défi technologique que ce projet s'efforce
de relever.

\subsection{Les limites du traitement manuel}

Le traitement manuel des commentaires étudiants présente plusieurs
limitations critiques, résumées dans le tableau \ref{tab:limites_manuel} :

\begin{table}[H]
\centering
\caption{Limitations du traitement manuel des commentaires}
\label{tab:limites_manuel}
\renewcommand{\arraystretch}{1.4}
\begin{tabular}{>{\bfseries}p{3.5cm} p{9cm}}
\toprule
\rowcolor{bleuprincipal!15}
\textbf{Limitation} & \textbf{Impact sur la qualité pédagogique} \\
\midrule
Lenteur & Délai de plusieurs semaines avant exploitation, rendant les insights obsolètes \\
\midrule
\rowcolor{gray!8}
Subjectivité & Interprétation variable selon l'analyste, biais cognitifs inconscients \\
\midrule
Coût élevé & Mobilisation de ressources humaines importantes pour une tâche répétitive \\
\midrule
\rowcolor{gray!8}
Non scalable & Impossible à maintenir avec l'augmentation des effectifs \\
\midrule
Exhaustivité & Tendance à se concentrer sur les cas extrêmes, ignorant les nuances \\
\bottomrule
\end{tabular}
\end{table}

% ──────────────────────────────────────────────────────────────────────────────
\section{Problématique Détaillée}
% ──────────────────────────────────────────────────────────────────────────────

La problématique centrale de ce projet se décline en plusieurs dimensions
complémentaires :

\subsection{Dimension linguistique}

Le français, avec ses 80 formes verbales conjuguées, ses nombreuses
exceptions morphologiques et ses structures syntaxiques complexes, représente
un défi particulier pour l'analyse automatique de sentiments. Les spécificités
suivantes méritent une attention particulière :

\begin{itemize}
  \item \textbf{La négation :} ``Ce cours n'est pas mauvais'' exprime en
        réalité quelque chose de positif, mais une approche naïve basée
        sur la présence du mot ``mauvais'' le classerait négativement.
  \item \textbf{L'intensification :} ``Vraiment excellent'', ``plutôt bien'',
        ``assez décevant'' — les adverbes modificateurs changent fondamentalement
        le niveau de sentiment exprimé.
  \item \textbf{L'ironie et le sarcasme :} ``Super, encore un cours
        incompréhensible...'' — nécessite une compréhension contextuelle
        profonde, impossible pour les approches lexicales.
\end{itemize}

\subsection{Dimension technologique}

Les approches traditionnelles d'analyse de sentiments basées sur des
dictionnaires de polarité (comme SentiWordNet ou VADER) ou des modèles
statistiques simples (SVM, Naive Bayes) présentent des performances
insuffisantes sur le français informel utilisé par les étudiants. Le
besoin d'un modèle profond pré-entraîné sur de larges corpus français
s'impose donc naturellement.

\subsection{Dimension applicative}

Au-delà du modèle seul, le système doit être intégré dans une chaîne
complète : ingestion des données, analyse, stockage, visualisation et
accès depuis des interfaces variées (web, mobile). Cette contrainte
d'intégration impose une architecture distribuée bien conçue.

% ──────────────────────────────────────────────────────────────────────────────
\section{Étude de l'Existant}
% ──────────────────────────────────────────────────────────────────────────────

\subsection{Solutions générales d'analyse de sentiments}

\subsubsection{TextBlob et VADER}

TextBlob et VADER sont des outils d'analyse de sentiments populaires pour
l'anglais. VADER utilise un dictionnaire de polarité augmenté de règles
grammaticales. Bien qu'efficaces pour l'anglais informel, ces outils offrent
des performances médiocres en français (précision inférieure à 60\% sur
des textes académiques).

\subsubsection{Google Cloud Natural Language API}

Google propose une API commerciale d'analyse de sentiments supportant
le français. Elle présente l'avantage d'être prête à l'emploi mais souffre
de plusieurs inconvénients majeurs dans notre contexte : coût par requête,
dépendance à une connexion internet, confidentialité des données étudiantes
et absence de personnalisation au vocabulaire académique.

\subsubsection{Modèles BERT multilingues}

\textbf{mBERT} (Multilingual BERT) et \textbf{XLM-RoBERTa} sont des modèles
Transformer entraînés sur des dizaines de langues simultanément. Bien que
supportant le français, leurs performances restent inférieures à un modèle
monolingue spécialisé, car ils partagent leur capacité entre 100+ langues.

\subsection{CamemBERT : l'état de l'art pour le français}

CamemBERT \cite{martin2020camembert} représente l'état de l'art pour le
Traitement du Langage Naturel en français. Développé conjointement par
Inria, Facebook AI Research et le CNRS, il a été entraîné sur 138 Go de
textes français provenant du corpus OSCAR \cite{ortiz2019oscar}. Ses
performances surpassent systématiquement mBERT sur toutes les tâches NLP
en français.

\begin{table}[H]
\centering
\caption{Comparaison des approches d'analyse de sentiments pour le français}
\label{tab:comparaison_approches}
\renewcommand{\arraystretch}{1.4}
\begin{tabular}{lcccp{3.5cm}}
\toprule
\rowcolor{bleuprincipal!15}
\textbf{Approche} & \textbf{Précision} & \textbf{Gratuit} &
\textbf{Hors-ligne} & \textbf{Personnalisable} \\
\midrule
VADER (anglais) & $<$55\% & Oui & Oui & Limitée \\
\rowcolor{gray!8}
TextBlob French & $\sim$58\% & Oui & Oui & Limitée \\
Google NL API & $\sim$72\% & Non & Non & Non \\
\rowcolor{gray!8}
mBERT & $\sim$70\% & Oui & Oui & Oui \\
XLM-RoBERTa & $\sim$74\% & Oui & Oui & Oui \\
\rowcolor{vertprincipal!15}
\textbf{CamemBERT (notre)} & \textbf{79,41\%} & \textbf{Oui} &
\textbf{Oui} & \textbf{Oui} \\
\bottomrule
\end{tabular}
\end{table}

% ──────────────────────────────────────────────────────────────────────────────
\section{Fondements Théoriques}
% ──────────────────────────────────────────────────────────────────────────────

\subsection{L'architecture Transformer}

Le Transformer \cite{vaswani2017attention} est une architecture de réseau
de neurones introduite en 2017 qui a révolutionné le domaine du NLP. Sa
caractéristique principale est le mécanisme d'\textbf{attention multi-têtes}
(\textit{Multi-Head Self-Attention}) qui permet au modèle de pondérer
l'importance de chaque mot dans une séquence par rapport à tous les autres
mots, capturant ainsi des dépendances à longue distance que les architectures
précédentes (RNN, LSTM) peinent à modéliser.

La formule d'attention est définie par :
\begin{equation}
\text{Attention}(Q, K, V) = \text{softmax}\left(\frac{QK^T}{\sqrt{d_k}}\right) V
\label{eq:attention}
\end{equation}

où $Q$ (Queries), $K$ (Keys) et $V$ (Values) sont des matrices de
projections linéaires, et $d_k$ est la dimension des clés.

\subsection{BERT et le pré-entraînement bidirectionnel}

\acrfull{bert} \cite{devlin2019bert} est un modèle Transformer qui exploite
le contexte bidirectionnel : contrairement aux modèles auto-régressifs comme
GPT qui lisent le texte de gauche à droite, BERT lit simultanément les
contextes gauche et droit de chaque token. Cette bidirectionnalité est
obtenue grâce à deux tâches de pré-entraînement :

\begin{itemize}
  \item \textbf{Masked Language Modeling (MLM) :} 15\% des tokens sont
        masqués aléatoirement, et le modèle doit les prédire.
  \item \textbf{Next Sentence Prediction (NSP) :} le modèle prédit si
        deux phrases sont consécutives dans le texte original.
\end{itemize}

\subsection{CamemBERT}

CamemBERT reprend l'architecture RoBERTa \cite{liu2019roberta} (optimisation
de BERT), mais est entraîné uniquement sur du texte français. Il utilise la
tokenisation \textbf{SentencePiece} \cite{kudo2018sentencepiece} avec un
vocabulaire de 32 000 tokens, permettant une meilleure représentation
des mots français et de leurs formes fléchies.

\subsection{Le fine-tuning par transfert d'apprentissage}

Le \textit{transfer learning} consiste à adapter un modèle pré-entraîné
sur une tâche générale à une tâche spécifique en le ré-entraînant sur un
petit jeu de données annoté. Pour la classification de sentiment, on ajoute
une couche de classification linéaire au-dessus du token \texttt{[CLS]}
de CamemBERT, puis on affine l'ensemble des paramètres :

\begin{equation}
P(\text{classe}_k | \mathbf{x}) = \text{softmax}(W \cdot \mathbf{h}_{CLS} + b)_k
\label{eq:classification}
\end{equation}

où $\mathbf{h}_{CLS}$ est la représentation du token \texttt{[CLS]},
$W$ est la matrice de poids de la couche de classification et $b$ le
biais.

% ──────────────────────────────────────────────────────────────────────────────
\section{Valeur Ajoutée du Projet}
% ──────────────────────────────────────────────────────────────────────────────

Notre système apporte plusieurs contributions concrètes par rapport aux
solutions existantes :

\begin{definitionbox}[Contributions principales]
\begin{enumerate}
  \item \textbf{Spécialisation domaine :} fine-tuning sur des commentaires
        académiques français, améliore la précision sur le domaine cible.
  \item \textbf{Système complet :} de la collecte de données à l'affichage
        mobile, en passant par l'API et le dashboard, le système est
        opérationnel bout en bout.
  \item \textbf{Temps réel :} analyse et retour instantané ($<$ 1 seconde),
        contrairement aux systèmes batch traditionnels.
  \item \textbf{Accessibilité :} application mobile Flutter pour les
        étudiants, dashboard web pour les enseignants.
  \item \textbf{Gratuité et confidentialité :} déployable entièrement
        en local, sans dépendance à des services cloud payants.
\end{enumerate}
\end{definitionbox}

\section*{Conclusion du Chapitre}

Ce chapitre a établi le contexte motivant notre projet, analysé les
limites des approches existantes et positionné CamemBERT comme le choix
technologique le plus pertinent pour l'analyse de sentiments en français
académique. Le chapitre suivant détaille l'analyse fonctionnelle et la
conception architecturale du système.

\cleardoublepage

% ==============================================================================
% CHAPITRE 2 : ANALYSE ET CONCEPTION
% ==============================================================================
\chapter{Analyse et Conception}

\begin{infobox}[Objectif du chapitre]
Ce chapitre présente l'analyse fonctionnelle et non fonctionnelle du système,
les diagrammes UML, l'architecture globale, le modèle de données et la
justification des choix technologiques.
\end{infobox}

% ──────────────────────────────────────────────────────────────────────────────
\section{Analyse des Besoins}
% ──────────────────────────────────────────────────────────────────────────────

\subsection{Identification des acteurs}

Le système interagit avec deux catégories d'acteurs principaux :

\begin{table}[H]
\centering
\caption{Acteurs du système}
\label{tab:acteurs}
\renewcommand{\arraystretch}{1.4}
\begin{tabular}{>{\bfseries}p{3cm} p{5cm} p{5cm}}
\toprule
\rowcolor{bleuprincipal!15}
\textbf{Acteur} & \textbf{Rôle} & \textbf{Interactions principales} \\
\midrule
Étudiant & Soumetteur de commentaires & Soumettre un commentaire, consulter le résultat, voir l'historique \\
\midrule
\rowcolor{gray!8}
Enseignant & Consommateur d'analyses & Consulter le dashboard, filtrer par matière, voir les statistiques \\
\midrule
Système IA & Acteur interne & Analyser le sentiment, retourner les scores \\
\bottomrule
\end{tabular}
\end{table}

\subsection{Besoins fonctionnels}

Les besoins fonctionnels représentent les fonctionnalités que le système
doit obligatoirement offrir :

\begin{enumerate}[label=\textbf{BF\arabic*}., leftmargin=3em]
  \item Le système doit permettre à un étudiant de soumettre un commentaire
        textuel d'au moins 5 caractères.
  \item Le système doit analyser automatiquement le sentiment du commentaire
        et retourner la classe prédite (Négatif/Neutre/Positif) avec les
        scores de confiance.
  \item Le système doit permettre d'associer un commentaire à une matière
        académique spécifique.
  \item Le système doit sauvegarder chaque commentaire analysé avec ses
        métadonnées dans la base de données.
  \item Le système doit permettre la consultation de l'historique des
        commentaires avec pagination.
  \item Le système doit offrir des filtres de recherche par sentiment,
        par matière et par date.
  \item Le système doit permettre la suppression d'un commentaire par
        son identifiant.
  \item Le système doit exposer des statistiques agrégées sur les analyses
        effectuées.
  \item L'application mobile doit permettre la navigation entre l'écran
        d'analyse et l'écran historique.
  \item Le tableau de bord doit afficher des visualisations interactives
        des tendances de sentiments.
\end{enumerate}

\subsection{Besoins non fonctionnels}

\begin{table}[H]
\centering
\caption{Besoins non fonctionnels}
\label{tab:bnf}
\renewcommand{\arraystretch}{1.4}
\begin{tabular}{>{\bfseries}p{3cm} p{10cm}}
\toprule
\rowcolor{bleuprincipal!15}
\textbf{Critère} & \textbf{Exigence} \\
\midrule
Performance & Le temps de réponse d'une analyse ne doit pas dépasser 2 secondes en conditions normales \\
\midrule
\rowcolor{gray!8}
Disponibilité & L'API doit être accessible 24h/24 une fois déployée \\
\midrule
Scalabilité & L'architecture doit permettre d'ajouter de nouvelles matières sans modification du code \\
\midrule
\rowcolor{gray!8}
Maintenabilité & Le code doit être modulaire, documenté et versionné via Git \\
\midrule
Sécurité & Les données étudiantes doivent être stockées localement, sans transfert vers des services tiers \\
\midrule
\rowcolor{gray!8}
Compatibilité & L'application mobile doit fonctionner sur Android 8.0+ \\
\midrule
Portabilité & Le système doit être déployable sous Linux Ubuntu sans dépendances propriétaires \\
\bottomrule
\end{tabular}
\end{table}

% ──────────────────────────────────────────────────────────────────────────────
\section{Diagrammes UML}
% ──────────────────────────────────────────────────────────────────────────────

\subsection{Diagramme de cas d'utilisation}

\begin{figure}[H]
\centering
\begin{tikzpicture}[
  actor/.style={circle, draw=bleuprincipal, fill=bleuprincipal!10,
                minimum size=1cm, font=\small\bfseries},
  usecase/.style={ellipse, draw=bleuprincipal, fill=white,
                  minimum width=3.5cm, minimum height=0.8cm,
                  font=\footnotesize, align=center, thick},
  system/.style={rectangle, draw=bleuprincipal!60, dashed,
                 minimum width=10cm, minimum height=8cm,
                 label={[font=\bfseries\small, text=bleuprincipal]
                        above:{Système d'Analyse de Sentiments}}}
]
  % Acteurs
  \node[actor] (etudiant) at (-5, 0) {};
  \node[below=0.1cm of etudiant, font=\small\bfseries] {Étudiant};

  \node[actor] (enseignant) at (7, 0) {};
  \node[below=0.1cm of enseignant, font=\small\bfseries] {Enseignant};

  % Frontière système
  \node[system] (sys) at (1, 0) {};

  % Cas d'utilisation
  \node[usecase] (uc1) at (1,  3.0) {Soumettre un\\commentaire};
  \node[usecase] (uc2) at (1,  1.8) {Voir le résultat\\d'analyse};
  \node[usecase] (uc3) at (1,  0.6) {Consulter\\l'historique};
  \node[usecase] (uc4) at (1, -0.6) {Filtrer par\\sentiment/matière};
  \node[usecase] (uc5) at (1, -1.8) {Supprimer un\\commentaire};
  \node[usecase] (uc6) at (1, -3.0) {Voir les\\statistiques};

  % Relations Étudiant
  \draw[->] (etudiant) -- (uc1);
  \draw[->] (etudiant) -- (uc2);
  \draw[->] (etudiant) -- (uc3);
  \draw[->] (etudiant) -- (uc4);
  \draw[->] (etudiant) -- (uc5);

  % Relations Enseignant
  \draw[->] (enseignant) -- (uc3);
  \draw[->] (enseignant) -- (uc4);
  \draw[->] (enseignant) -- (uc6);
  \draw[->] (enseignant) -- (uc5);

\end{tikzpicture}
\caption{Diagramme de cas d'utilisation du système}
\label{fig:uml_usecase}
\end{figure}

\subsection{Diagramme de séquence — Analyse d'un commentaire}

\begin{figure}[H]
\centering
\begin{tikzpicture}[
  lifeline/.style={draw=bleuprincipal!50, dashed, thick},
  box/.style={rectangle, draw=bleuprincipal, fill=bleuprincipal!10,
              minimum width=2.2cm, minimum height=0.7cm,
              font=\small\bfseries, align=center},
  msg/.style={->, thick, color=bleuprincipal!80},
  ret/.style={->, thick, dashed, color=vertprincipal!80}
]
  % En-têtes
  \node[box] (user)    at (0,0)   {Étudiant\\(Flutter)};
  \node[box] (api)     at (4,0)   {FastAPI\\Backend};
  \node[box] (service) at (8,0)   {Service\\Prédiction};
  \node[box] (model)   at (11.5,0){CamemBERT\\Model};
  \node[box] (db)      at (15,0)  {PostgreSQL};

  % Lignes de vie
  \draw[lifeline] (0,-0.35) -- (0,-10);
  \draw[lifeline] (4,-0.35) -- (4,-10);
  \draw[lifeline] (8,-0.35) -- (8,-10);
  \draw[lifeline] (11.5,-0.35) -- (11.5,-10);
  \draw[lifeline] (15,-0.35) -- (15,-10);

  % Messages
  \draw[msg] (0,-1.2) -- node[above,font=\footnotesize]{POST /api/predict} (4,-1.2);
  \draw[msg] (4,-2.0) -- node[above,font=\footnotesize]{predict(texte)} (8,-2.0);
  \draw[msg] (8,-2.8) -- node[above,font=\footnotesize]{nettoyer(texte)} (8,-2.8);
  \draw[msg] (8,-3.5) -- node[above,font=\footnotesize]{tokeniser(texte)} (11.5,-3.5);
  \draw[msg] (11.5,-4.3) -- node[above,font=\footnotesize]{forward()} (11.5,-4.3);
  \draw[ret] (11.5,-5.1) -- node[above,font=\footnotesize]{logits} (8,-5.1);
  \draw[msg] (8,-5.9) -- node[above,font=\footnotesize]{softmax(logits)} (8,-5.9);
  \draw[ret] (8,-6.7) -- node[above,font=\footnotesize]{\{sentiment, scores\}} (4,-6.7);
  \draw[msg] (4,-7.5) -- node[above,font=\footnotesize]{INSERT commentaire} (15,-7.5);
  \draw[ret] (15,-8.3) -- node[above,font=\footnotesize]{id, date\_creation} (4,-8.3);
  \draw[ret] (4,-9.1) -- node[above,font=\footnotesize]{JSON résultat complet} (0,-9.1);

\end{tikzpicture}
\caption{Diagramme de séquence : analyse d'un commentaire}
\label{fig:seq_analyse}
\end{figure}

\subsection{Diagramme de classes}

\begin{figure}[H]
\centering
\begin{tikzpicture}[
  class/.style={rectangle, draw=bleuprincipal, fill=white,
                minimum width=5cm, thick},
  attr/.style={font=\small\ttfamily},
  meth/.style={font=\small\ttfamily, color=bleuprincipal}
]

% Classe Commentaire (DB)
\node[class] (commentaire) at (0,0) {
  \begin{tabular}{l}
    \rowcolor{bleuprincipal!20}
    \multicolumn{1}{c}{\textbf{Commentaire (DB)}} \\
    \hline
    \small\ttfamily id : Integer (PK) \\
    \small\ttfamily texte\_original : Text \\
    \small\ttfamily texte\_nettoye : Text \\
    \small\ttfamily matiere : String \\
    \small\ttfamily sentiment\_predit : Integer \\
    \small\ttfamily label\_sentiment : String \\
    \small\ttfamily score\_negatif : Float \\
    \small\ttfamily score\_neutre : Float \\
    \small\ttfamily score\_positif : Float \\
    \small\ttfamily date\_creation : DateTime \\
    \hline
    \small\ttfamily to\_dict() : dict \\
  \end{tabular}
};

% Classe ServicePrediction
\node[class] (service) at (8.5, 2) {
  \begin{tabular}{l}
    \rowcolor{vertprincipal!20}
    \multicolumn{1}{c}{\textbf{ServicePrediction}} \\
    \hline
    \small\ttfamily \_instance : Singleton \\
    \small\ttfamily modele : CamemBERT \\
    \small\ttfamily tokeniseur : Tokenizer \\
    \small\ttfamily device : torch.device \\
    \hline
    \small\ttfamily charger() : void \\
    \small\ttfamily predict(texte) : dict \\
  \end{tabular}
};

% Classe Commentaire (Flutter)
\node[class] (flutter) at (8.5, -2.5) {
  \begin{tabular}{l}
    \rowcolor{orangeaccent!15}
    \multicolumn{1}{c}{\textbf{Commentaire (Flutter)}} \\
    \hline
    \small\ttfamily id : int \\
    \small\ttfamily texteOriginal : String \\
    \small\ttfamily matiere : String? \\
    \small\ttfamily sentimentPredit : int \\
    \small\ttfamily scores : Map \\
    \hline
    \small\ttfamily fromJson(json) \\
    \small\ttfamily emoji : String \\
    \small\ttfamily scoreConfiance : double \\
  \end{tabular}
};

% Flèches
\draw[->, thick, color=bleuprincipal] (service) -- node[above, font=\small]{utilise} (commentaire);
\draw[->, thick, dashed, color=orangeaccent] (flutter) -- node[right, font=\small]{mappe} (commentaire);

\end{tikzpicture}
\caption{Diagramme de classes simplifié}
\label{fig:class_diagram}
\end{figure}

% ──────────────────────────────────────────────────────────────────────────────
\section{Architecture Globale du Système}
% ──────────────────────────────────────────────────────────────────────────────

\subsection{Vue d'ensemble}

\begin{figure}[H]
\centering
\begin{tikzpicture}[
  layer/.style={rectangle, rounded corners=6pt, minimum width=3.5cm,
                minimum height=1.2cm, align=center, font=\small\bfseries,
                thick, draw},
  arrow/.style={->, thick, color=bleuprincipal!70},
  label/.style={font=\tiny, color=gray}
]

% Couche présentation
\node[layer, fill=bleuclaire!20, draw=bleuclaire] (flutter) at (0, 6)
  {\faIcon{mobile-alt} Application\\Flutter};
\node[layer, fill=bleuclaire!20, draw=bleuclaire] (streamlit) at (5, 6)
  {\faIcon{chart-bar} Dashboard\\Streamlit};

% Couche API
\node[layer, fill=bleuprincipal!15, draw=bleuprincipal, minimum width=9cm]
  (fastapi) at (2.5, 4) {\faIcon{server} API REST — FastAPI (port 8000)};

% Couche services
\node[layer, fill=vertprincipal!15, draw=vertprincipal] (predict) at (0, 2)
  {\faIcon{brain} Service\\Prédiction};
\node[layer, fill=vertprincipal!15, draw=vertprincipal] (router) at (5, 2)
  {\faIcon{route} Routeur\\Commentaires};

% Couche IA
\node[layer, fill=orangeaccent!15, draw=orangeaccent] (camembert) at (0, 0)
  {\faIcon{robot} CamemBERT\\Fine-tuné};

% Couche données
\node[layer, fill=rougeprincipal!10, draw=rougeprincipal] (postgres) at (5, 0)
  {\faIcon{database} PostgreSQL\\sentiment\_etudiants};

% Connexions
\draw[arrow] (flutter) -- node[label, right]{HTTP/JSON} (fastapi);
\draw[arrow] (streamlit) -- node[label, left]{HTTP/JSON} (fastapi);
\draw[arrow] (fastapi) -- (predict);
\draw[arrow] (fastapi) -- (router);
\draw[arrow] (predict) -- (camembert);
\draw[arrow] (router) -- node[label, right]{SQLAlchemy ORM} (postgres);
\draw[arrow] (predict) -- node[label, below]{sauvegarder} (router);

% Légende
\node[font=\tiny, gray] at (2.5, -1.2)
  {Architecture 4 couches : Présentation — API — Métier — Données};

\end{tikzpicture}
\caption{Architecture globale du système en 4 couches}
\label{fig:architecture}
\end{figure}

\subsection{Description des couches}

\subsubsection{Couche Présentation}
Comprend deux interfaces utilisateur distinctes selon le profil d'utilisateur :
l'application mobile Flutter pour les étudiants, et le dashboard Streamlit
pour les enseignants/administrateurs. Les deux communiquent exclusivement
via l'API REST.

\subsubsection{Couche API (FastAPI)}
Le point d'entrée unique du système. FastAPI gère la validation des
requêtes entrantes via les schémas Pydantic, l'authentification CORS,
le routage vers les services appropriés et la sérialisation des réponses
JSON.

\subsubsection{Couche Métier}
Contient la logique fonctionnelle : le \texttt{ServicePrediction} (singleton
gérant le modèle CamemBERT) et le routeur de commentaires gérant les
opérations CRUD.

\subsubsection{Couche Données}
PostgreSQL stocke toutes les analyses effectuées. SQLAlchemy ORM fait
office de couche d'abstraction, permettant d'interagir avec la base via
des objets Python sans écrire de SQL brut.

% ──────────────────────────────────────────────────────────────────────────────
\section{Modèle de Données}
% ──────────────────────────────────────────────────────────────────────────────

\subsection{Modèle Conceptuel de Données (MCD)}

Le modèle de données est intentionnellement simple : une seule entité
principale \textbf{Commentaire} regroupe toutes les informations relatives
à une analyse de sentiment.

\begin{figure}[H]
\centering
\begin{tikzpicture}[
  entity/.style={rectangle, draw=bleuprincipal, fill=bleuprincipal!10,
                 minimum width=6cm, font=\small\bfseries},
  attr/.style={font=\footnotesize\ttfamily}
]
\node[entity] (header) at (0,0) {COMMENTAIRE};
\node[rectangle, draw=bleuprincipal!60, fill=white, minimum width=6cm,
      align=left, font=\footnotesize\ttfamily] (body) at (0,-3.5) {
  \begin{tabular}{ll}
    \textbf{PK} & id : SERIAL \\
    & texte\_original : TEXT \\
    & texte\_nettoye : TEXT \\
    & matiere : VARCHAR(100) \\
    & sentiment\_predit : INTEGER \\
    & label\_sentiment : VARCHAR(20) \\
    & score\_negatif : FLOAT \\
    & score\_neutre : FLOAT \\
    & score\_positif : FLOAT \\
    & date\_creation : TIMESTAMP \\
  \end{tabular}
};
\draw[bleuprincipal] (header.south west) -- (body.north west);
\draw[bleuprincipal] (header.south east) -- (body.north east);
\end{tikzpicture}
\caption{Modèle Conceptuel de Données}
\label{fig:mcd}
\end{figure}

\subsection{Modèle Logique de Données (MLD)}

Le MLD correspond directement à la table SQL créée par SQLAlchemy :

\begin{lstlisting}[style=styleSQL, caption={Schéma SQL de la table commentaires}, label={lst:sql_schema}]
CREATE TABLE commentaires (
    id               SERIAL          PRIMARY KEY,
    texte_original   TEXT            NOT NULL,
    texte_nettoye    TEXT,
    matiere          VARCHAR(100),
    sentiment_predit INTEGER         NOT NULL,  -- 0, 1 ou 2
    label_sentiment  VARCHAR(20)     NOT NULL,  -- 'Negatif', 'Neutre', 'Positif'
    score_negatif    DOUBLE PRECISION,
    score_neutre     DOUBLE PRECISION,
    score_positif    DOUBLE PRECISION,
    date_creation    TIMESTAMP       DEFAULT NOW()
);
\end{lstlisting}

% ──────────────────────────────────────────────────────────────────────────────
\section{Justification des Choix Technologiques}
% ──────────────────────────────────────────────────────────────────────────────

\begin{table}[H]
\centering
\caption{Justification des choix technologiques}
\label{tab:choix_tech}
\renewcommand{\arraystretch}{1.5}
\begin{tabular}{p{2.5cm} p{3cm} p{7cm}}
\toprule
\rowcolor{bleuprincipal!15}
\textbf{Technologie} & \textbf{Rôle} & \textbf{Justification} \\
\midrule
CamemBERT & Modèle NLP & État de l'art pour le français, pré-entraîné sur 138 Go, surpasse tous les modèles multilingues \\
\midrule
\rowcolor{gray!8}
PyTorch & Framework DL & Standard industriel, support natif HuggingFace, flexible et débuggable \\
\midrule
HuggingFace & Bibliothèque & Simplifie le chargement et fine-tuning des Transformers, vaste communauté \\
\midrule
\rowcolor{gray!8}
FastAPI & API REST & Le plus rapide des frameworks Python API, typage fort, documentation automatique Swagger \\
\midrule
SQLAlchemy & ORM & Abstraction DB portable, sécurité SQL injection, migrations faciles \\
\midrule
\rowcolor{gray!8}
PostgreSQL & Base de données & ACID, performances, support JSON natif, open-source et robuste \\
\midrule
Streamlit & Dashboard & Python pur, déploiement instantané, intégration Plotly native \\
\midrule
\rowcolor{gray!8}
Flutter & Application mobile & Cross-platform (Android/iOS/Web), Material 3, performances natives \\
\midrule
Dart & Langage mobile & Typage fort, compilation AOT, excellent pour les UIs réactives \\
\bottomrule
\end{tabular}
\end{table}

\section*{Conclusion du Chapitre}

L'analyse fonctionnelle a permis d'identifier précisément les besoins des
utilisateurs et de concevoir une architecture répondant à ces besoins de
manière optimale. Le choix de CamemBERT comme modèle central, combiné à
FastAPI, PostgreSQL et Flutter, forme un écosystème cohérent et moderne.
Le chapitre suivant détaille l'implémentation concrète de chaque composant.

\cleardoublepage

% ==============================================================================
% CHAPITRE 3 : IMPLÉMENTATION
% ==============================================================================
\chapter{Implémentation}

\begin{infobox}[Objectif du chapitre]
Ce chapitre décrit en détail l'environnement de développement, la structure
du projet, les algorithmes clés et les parties importantes du code source
de chaque composant du système.
\end{infobox}

% ──────────────────────────────────────────────────────────────────────────────
\section{Environnements de Développement}
% ──────────────────────────────────────────────────────────────────────────────

\subsection{Environnement matériel}

\begin{table}[H]
\centering
\caption{Configuration matérielle de développement}
\label{tab:hardware}
\renewcommand{\arraystretch}{1.4}
\begin{tabular}{lll}
\toprule
\rowcolor{bleuprincipal!15}
\textbf{Composant} & \textbf{Spécification} & \textbf{Impact} \\
\midrule
Processeur & Intel/AMD x86-64 & Exécution du modèle sur CPU \\
\rowcolor{gray!8}
Mémoire RAM & 8 Go minimum & Chargement du modèle ($\sim$430 Mo) \\
Stockage & 20 Go disponibles & Modèle + SDK Android + venv \\
\rowcolor{gray!8}
OS & Ubuntu 24.04 LTS & Environnement de production Linux \\
Réseau & Wi-Fi / Ethernet & Communication émulateur $\leftrightarrow$ API \\
\bottomrule
\end{tabular}
\end{table}

\subsection{Environnement logiciel}

\begin{table}[H]
\centering
\caption{Logiciels et versions utilisés}
\label{tab:software}
\renewcommand{\arraystretch}{1.4}
\begin{tabular}{lll}
\toprule
\rowcolor{bleuprincipal!15}
\textbf{Logiciel} & \textbf{Version} & \textbf{Rôle} \\
\midrule
Python & 3.12 & Langage backend et IA \\
\rowcolor{gray!8}
PyTorch & 2.x & Framework Deep Learning \\
Transformers (HuggingFace) & 4.x & Modèles pré-entraînés \\
\rowcolor{gray!8}
FastAPI & 0.110+ & Framework API REST \\
SQLAlchemy & 2.0+ & ORM Python-PostgreSQL \\
\rowcolor{gray!8}
PostgreSQL & 16 & Base de données relationnelle \\
Streamlit & 1.x & Dashboard web Python \\
\rowcolor{gray!8}
Flutter SDK & 3.24.5 & Développement application mobile \\
Dart & 3.x & Langage de l'application mobile \\
\rowcolor{gray!8}
Android Studio & 2024.2 & Émulateur Android, IDE Flutter \\
VS Code & 1.119 & Éditeur de code principal \\
\rowcolor{gray!8}
Git & 2.x & Contrôle de version \\
\bottomrule
\end{tabular}
\end{table}

% ──────────────────────────────────────────────────────────────────────────────
\section{Structure du Projet}
% ──────────────────────────────────────────────────────────────────────────────

\begin{lstlisting}[
  language=bash, caption={Structure complète du projet},
  label={lst:structure}, basicstyle=\ttfamily\footnotesize,
  frame=single, backgroundcolor=\color{griscode},
  numbers=none, xleftmargin=0.5cm
]
sentiment_etudiants/
|-- ai_model/                  # Module IA (CamemBERT)
|   |-- preprocessing/
|   |   |-- nettoyage_texte.py # Pipeline de nettoyage NLP
|   |   |-- tokenisation.py    # Tokenisation CamemBERT
|   |   `-- preparer_donnees.py# Pipeline preparation dataset
|   |-- modele_camembert.py    # Chargement/sauvegarde du modele
|   |-- entrainement.py        # Script d'entrainement (fine-tuning)
|   |-- evaluation.py          # Metriques: accuracy, F1, matrice
|   `-- saved_model/           # Modele fine-tune sauvegarde
|-- backend/                   # API FastAPI
|   |-- database.py            # Connexion PostgreSQL via SQLAlchemy
|   |-- models/modeles_db.py   # ORM: table commentaires
|   |-- routers/commentaires.py# Routes REST: predict, CRUD
|   |-- services/prediction_service.py # Singleton inference
|   `-- main.py                # Point d'entree FastAPI
|-- dashboard/                 # Tableau de bord Streamlit
|   |-- app.py                 # Page d'accueil
|   |-- components/api_client.py# Client HTTP vers l'API
|   `-- pages/                 # Pages Streamlit multi-pages
|       |-- 1_Analyse.py       # Analyse individuelle
|       |-- 2_Statistiques.py  # Graphiques Plotly
|       `-- 3_Historique.py    # Liste et filtres
|-- mobile_app/                # Application Flutter
|   `-- lib/
|       |-- main.dart          # Point d'entree + navigation
|       |-- models/commentaire.dart  # Modele de donnees Dart
|       |-- services/api_service.dart# Client HTTP Dart
|       |-- widgets/sentiment_card.dart # Widget reutilisable
|       `-- screens/
|           |-- analyse_screen.dart  # Ecran d'analyse
|           `-- historique_screen.dart# Ecran historique
`-- datasets/
    |-- raw/commentaires_etudiants.csv # 225 commentaires
    `-- processed/{train,validation,test}.csv
\end{lstlisting}

% ──────────────────────────────────────────────────────────────────────────────
\section{Module IA — Pipeline CamemBERT}
% ──────────────────────────────────────────────────────────────────────────────

\subsection{Prétraitement du texte}

Avant toute tokenisation, les commentaires bruts sont nettoyés par un
pipeline séquentiel défini dans \texttt{nettoyage\_texte.py} :

\begin{lstlisting}[style=stylePython, caption={Pipeline de nettoyage NLP}, label={lst:nettoyage}]
def nettoyer_commentaire(texte: str) -> str:
    """Pipeline complet de nettoyage d'un commentaire."""
    if not isinstance(texte, str) or len(texte.strip()) == 0:
        return ""
    # Etape 1 : Normalisation Unicode (NFC)
    texte = unicodedata.normalize("NFC", texte)
    texte = texte.replace(u"«", '"').replace(u"»", '"')
    # Etape 2 : Suppression des URLs, emails et caracteres speciaux
    texte = re.sub(r"http\S+|www\S+", "", texte)
    texte = re.sub(r"\S+@\S+", "", texte)
    texte = re.sub(r"[^\w\sA-Za-zÀ-ž'',;:!?.\-]", " ", texte)
    # Etape 3 : Mise en minuscules
    texte = texte.lower()
    # Etape 4 : Normalisation des espaces
    texte = re.sub(r"\s+", " ", texte).strip()
    return texte
\end{lstlisting}

Ce pipeline traite les cas suivants :
\begin{itemize}
  \item \textbf{Normalisation Unicode :} harmonise les caractères composés
        (é peut être représenté de deux façons en Unicode).
  \item \textbf{Guillemets typographiques :} « » sont remplacés par des
        guillemets droits ASCII.
  \item \textbf{URLs et emails :} supprimés car non pertinents pour l'analyse
        de sentiment.
  \item \textbf{Minuscules :} CamemBERT est sensible à la casse, la mise
        en minuscules améliore la cohérence.
\end{itemize}

\subsection{Jeu de données}

Le jeu de données a été constitué manuellement avec \textbf{225 commentaires
étudiants} en français, répartis équitablement en 3 classes :

\begin{table}[H]
\centering
\caption{Répartition du jeu de données}
\label{tab:dataset}
\renewcommand{\arraystretch}{1.4}
\begin{tabular}{lccc}
\toprule
\rowcolor{bleuprincipal!15}
\textbf{Classe} & \textbf{Étiquette} & \textbf{Exemples} & \textbf{Proportion} \\
\midrule
Négatif & 0 & 75 & 33,33\% \\
\rowcolor{gray!8}
Neutre & 1 & 75 & 33,33\% \\
Positif & 2 & 75 & 33,33\% \\
\midrule
\rowcolor{bleuprincipal!8}
\textbf{Total} & -- & \textbf{225} & \textbf{100\%} \\
\bottomrule
\end{tabular}
\end{table}

Le découpage train/validation/test respecte une stratification :

\begin{table}[H]
\centering
\caption{Découpage du dataset (stratifié)}
\label{tab:split}
\renewcommand{\arraystretch}{1.4}
\begin{tabular}{lcccc}
\toprule
\rowcolor{bleuprincipal!15}
\textbf{Ensemble} & \textbf{Taille} & \textbf{Proportion} &
\textbf{Usage} & \textbf{Batches} \\
\midrule
Entraînement & 157 & $\sim$70\% & Mise à jour des poids & 10 \\
\rowcolor{gray!8}
Validation & 34 & $\sim$15\% & Sélection du meilleur modèle & 3 \\
Test & 34 & $\sim$15\% & Évaluation finale non biaisée & 3 \\
\bottomrule
\end{tabular}
\end{table}

\subsection{Architecture du modèle fine-tuné}

Le modèle CamemBERT fine-tuné est composé de :

\begin{enumerate}
  \item \textbf{Couche d'entrée :} tokenisation SentencePiece (vocabulaire
        32 000 tokens, séquences de max 128 tokens).
  \item \textbf{12 blocs Transformer :} chaque bloc contient un mécanisme
        d'attention multi-têtes (12 têtes, dimension 768) et un réseau
        Feed-Forward (dimension 3072).
  \item \textbf{Couche de classification :} projection linéaire
        $768 \rightarrow 3$ suivie d'un Softmax.
\end{enumerate}

Le nombre total de paramètres est de \textbf{110,6 millions}, dont les
poids de la couche de classification ($768 \times 3 + 3 = 2307$ paramètres
ajoutés).

\subsection{Boucle d'entraînement}

\begin{lstlisting}[style=stylePython, caption={Boucle d'entraînement fine-tuning CamemBERT}, label={lst:training}]
def entrainer_une_epoch(modele, loader_train, optimiseur, scheduler,
                        device, num_epoch):
    modele.train()
    perte_totale = 0.0
    nb_correct, nb_total = 0, 0

    for i, batch in enumerate(loader_train):
        input_ids      = batch["input_ids"].to(device)
        attention_mask = batch["attention_mask"].to(device)
        labels         = batch["labels"].to(device)

        # 1. Remise a zero des gradients
        optimiseur.zero_grad()

        # 2. Passage avant : calcul de la perte cross-entropique
        sorties = modele(input_ids=input_ids,
                         attention_mask=attention_mask,
                         labels=labels)
        perte = sorties.loss

        # 3. Retropropagation : calcul des gradients
        perte.backward()

        # 4. Ecrêtage des gradients (previent l'explosion)
        torch.nn.utils.clip_grad_norm_(modele.parameters(), max_norm=1.0)

        # 5. Mise a jour des poids
        optimiseur.step()
        scheduler.step()  # Ajuste le learning rate

        perte_totale += perte.item()
        predictions = torch.argmax(sorties.logits, dim=1)
        nb_correct  += (predictions == labels).sum().item()
        nb_total    += labels.size(0)

    return perte_totale / len(loader_train), nb_correct / nb_total * 100
\end{lstlisting}

\textbf{Points techniques importants :}
\begin{itemize}
  \item \texttt{optimiseur.zero\_grad()} : réinitialise les gradients
        accumulés du batch précédent. Sans cela, les gradients s'accumulent
        à chaque batch, causant des mises à jour erronées.
  \item \texttt{clip\_grad\_norm\_(max\_norm=1.0)} : technique standard pour
        les Transformers évitant l'explosion des gradients (gradient clipping).
  \item \texttt{scheduler.step()} : ajuste le taux d'apprentissage selon
        le scheduler linéaire avec warmup défini.
\end{itemize}

\subsection{Hyperparamètres d'entraînement}

\begin{table}[H]
\centering
\caption{Hyperparamètres du fine-tuning}
\label{tab:hyperparams}
\renewcommand{\arraystretch}{1.4}
\begin{tabular}{lll}
\toprule
\rowcolor{bleuprincipal!15}
\textbf{Hyperparamètre} & \textbf{Valeur} & \textbf{Justification} \\
\midrule
Taux d'apprentissage & $2 \times 10^{-5}$ & Standard fine-tuning BERT, évite le catastrophic forgetting \\
\rowcolor{gray!8}
Taille de lot (batch) & 16 & Compromis mémoire/variance du gradient \\
Époques & 3 & Suffisant pour un petit dataset, évite le surapprentissage \\
\rowcolor{gray!8}
Max tokens & 128 & Couvre 95\%+ des commentaires, économise la mémoire \\
Optimiseur & AdamW & Standard pour les Transformers, weight decay régularise \\
\rowcolor{gray!8}
Warmup ratio & 10\% & Stabilise le début de l'entraînement \\
Gradient clipping & 1.0 & Prévient l'explosion des gradients \\
\rowcolor{gray!8}
Weight decay & 0.01 & Régularisation L2 pour réduire le surapprentissage \\
\bottomrule
\end{tabular}
\end{table}

% ──────────────────────────────────────────────────────────────────────────────
\section{Module Backend — FastAPI}
% ──────────────────────────────────────────────────────────────────────────────

\subsection{Démarrage du serveur et lifespan}

\begin{lstlisting}[style=stylePython, caption={Initialisation FastAPI avec lifespan}, label={lst:fastapi_main}]
@asynccontextmanager
async def lifespan(app: FastAPI):
    """Initialisation au demarrage, nettoyage a l'arret."""
    print("[DEMARRAGE] Initialisation de l'API...")
    creer_tables()           # Creation des tables PostgreSQL
    service_prediction.charger()  # Chargement du modele CamemBERT en RAM
    print("[DEMARRAGE] API prete.")
    yield                    # Le serveur tourne ici
    print("[ARRET] Fermeture de l'API.")

app = FastAPI(
    title="API Analyse de Sentiments Etudiants",
    version="1.0.0",
    lifespan=lifespan
)

app.add_middleware(CORSMiddleware,
    allow_origins=["*"],  # Flutter + Streamlit peuvent appeler l'API
    allow_methods=["*"],
    allow_headers=["*"]
)
\end{lstlisting}

Le pattern \texttt{lifespan} remplace les anciens \texttt{startup\_event}
et \texttt{shutdown\_event} de FastAPI. Il garantit que le modèle CamemBERT
est chargé une seule fois en mémoire au démarrage du serveur, évitant un
rechargement coûteux à chaque requête.

\subsection{Service de prédiction — Pattern Singleton}

\begin{lstlisting}[style=stylePython, caption={Implémentation du pattern Singleton}, label={lst:singleton}]
class ServicePrediction:
    """
    Singleton : une seule instance partagee dans toute l'application.
    Garantit que le modele CamemBERT (430 Mo) n'est charge qu'une fois.
    """
    _instance = None

    def __new__(cls):
        if cls._instance is None:
            cls._instance = super().__new__(cls)
            cls._instance._initialise = False
        return cls._instance

    def predict(self, texte: str) -> dict:
        texte_nettoye = nettoyer_commentaire(texte)

        encodage = self.tokeniseur(
            texte_nettoye,
            max_length=128,
            padding="max_length",
            truncation=True,
            return_tensors="pt"
        )

        with torch.no_grad():  # Desactive le calcul des gradients
            sorties = self.modele(
                input_ids=encodage["input_ids"].to(self.device),
                attention_mask=encodage["attention_mask"].to(self.device)
            )

        # Conversion logits -> probabilites
        probs = F.softmax(sorties.logits, dim=1).squeeze().tolist()
        classe = int(torch.argmax(sorties.logits, dim=1).item())

        return {
            "sentiment_predit": classe,
            "label_sentiment" : LABELS[classe],
            "score_negatif"   : round(probs[0], 4),
            "score_neutre"    : round(probs[1], 4),
            "score_positif"   : round(probs[2], 4),
        }
\end{lstlisting}

\textbf{Pourquoi \texttt{torch.no\_grad()} ?} Lors de l'inférence (prédiction),
nous n'avons pas besoin de calculer les gradients (ça, c'est pour
l'entraînement). Cette directive réduit la consommation mémoire de 50\%
et accélère les calculs.

\subsection{Validation des données avec Pydantic}

\begin{lstlisting}[style=stylePython, caption={Schémas de validation Pydantic}, label={lst:pydantic}]
class RequetePrediction(BaseModel):
    """Corps de la requete POST /api/predict"""
    texte   : str           = Field(..., min_length=5, max_length=2000,
                              description="Commentaire a analyser")
    matiere : Optional[str] = Field(None, max_length=100)

class ReponseCommentaire(BaseModel):
    """Schema de reponse standardise"""
    id               : int
    texte_original   : str
    matiere          : Optional[str]
    sentiment_predit : int
    label_sentiment  : str
    scores           : dict
    date_creation    : Optional[str]

    class Config:
        from_attributes = True  # Lecture depuis un objet SQLAlchemy
\end{lstlisting}

Pydantic garantit que :
\begin{itemize}
  \item Le texte fait au moins 5 caractères (sinon HTTP 422).
  \item Les types sont corrects (int, str, etc.).
  \item Les champs optionnels sont traités correctement.
  \item La sérialisation JSON est automatique.
\end{itemize}

\subsection{Routes API — Récapitulatif}

\begin{table}[H]
\centering
\caption{Endpoints de l'API REST}
\label{tab:endpoints}
\renewcommand{\arraystretch}{1.5}
\begin{tabular}{llp{3.5cm}p{3.5cm}}
\toprule
\rowcolor{bleuprincipal!15}
\textbf{Méthode} & \textbf{Route} & \textbf{Paramètres} & \textbf{Description} \\
\midrule
\rowcolor{vertprincipal!10}
POST & \texttt{/api/predict} & \texttt{\{texte, matiere?\}} & Analyse + sauvegarde \\
GET & \texttt{/api/commentaires} & \texttt{?limite, sentiment, matiere} & Liste paginée filtrée \\
\rowcolor{gray!8}
GET & \texttt{/api/commentaires/\{id\}} & \texttt{id} (chemin) & Détail d'un commentaire \\
GET & \texttt{/api/stats} & -- & Statistiques agrégées \\
\rowcolor{rougeprincipal!10}
DELETE & \texttt{/api/commentaires/\{id\}} & \texttt{id} (chemin) & Suppression \\
GET & \texttt{/health} & -- & Statut du serveur \\
\bottomrule
\end{tabular}
\end{table}

% ──────────────────────────────────────────────────────────────────────────────
\section{Module Application Mobile — Flutter}
% ──────────────────────────────────────────────────────────────────────────────

\subsection{Service API — Communication HTTP}

\begin{lstlisting}[style=styleDart, caption={Service HTTP Flutter}, label={lst:flutter_api}]
class ApiService {
  static const String _baseUrl = 'http://10.0.2.2:8000/api';
  static const Duration _timeout = Duration(seconds: 30);

  static Future<Commentaire?> analyserCommentaire({
    required String texte,
    String? matiere,
  }) async {
    try {
      final body = <String, dynamic>{'texte': texte};
      if (matiere != null) body['matiere'] = matiere;

      final reponse = await http.post(
        Uri.parse('$_baseUrl/predict'),
        headers: {'Content-Type': 'application/json; charset=utf-8'},
        body: jsonEncode(body),
      ).timeout(_timeout);

      if (reponse.statusCode == 200) {
        final json = jsonDecode(utf8.decode(reponse.bodyBytes));
        return Commentaire.fromJson(json);
      }
      return null;
    } catch (_) {
      return null;  // Gestion silencieuse des erreurs reseau
    }
  }
}
\end{lstlisting}

\textbf{Pourquoi \texttt{10.0.2.2} ?} L'émulateur Android est une machine
virtuelle qui possède sa propre interface réseau. Pour accéder au
\texttt{localhost} de l'ordinateur hôte depuis l'émulateur, Android
redirige l'adresse \texttt{10.0.2.2} vers \texttt{127.0.0.1} de la machine
hôte. Sur un vrai téléphone, il faudrait utiliser l'IP locale du PC.

\subsection{Navigation par onglets avec IndexedStack}

\begin{lstlisting}[style=styleDart, caption={Navigation IndexedStack}, label={lst:navigation}]
class _AccueilPageState extends State<AccueilPage> {
  int _ongletActif = 0;

  final List<Widget> _ecrans = const [
    AnalyseScreen(),     // Index 0
    HistoriqueScreen(),  // Index 1
  ];

  @override
  Widget build(BuildContext context) {
    return Scaffold(
      body: IndexedStack(
        index: _ongletActif,
        children: _ecrans,
      ),
      bottomNavigationBar: NavigationBar(
        selectedIndex: _ongletActif,
        onDestinationSelected: (i) => setState(() => _ongletActif = i),
        // ...
      ),
    );
  }
}
\end{lstlisting}

\texttt{IndexedStack} maintient les deux écrans en vie simultanément en
mémoire. L'écran visible est déterminé par l'index. Ce choix architectural
préserve l'état de chaque écran lors de la navigation (la liste de l'historique
n'est pas rechargée à chaque changement d'onglet).

% ──────────────────────────────────────────────────────────────────────────────
\section{Gestion des Erreurs et Performances}
% ──────────────────────────────────────────────────────────────────────────────

\subsection{Gestion des erreurs}

\begin{table}[H]
\centering
\caption{Stratégie de gestion des erreurs}
\label{tab:errors}
\renewcommand{\arraystretch}{1.4}
\begin{tabular}{p{3cm} p{3.5cm} p{6cm}}
\toprule
\rowcolor{bleuprincipal!15}
\textbf{Couche} & \textbf{Erreur possible} & \textbf{Gestion} \\
\midrule
API FastAPI & Texte trop court & HTTP 422 (Unprocessable Entity) + message \\
\rowcolor{gray!8}
API FastAPI & Erreur modèle & HTTP 500 + message d'erreur technique \\
API FastAPI & ID inexistant & HTTP 404 (Not Found) \\
\rowcolor{gray!8}
Flutter & Timeout réseau & \texttt{null} retourné, message affiché \\
Flutter & Serveur injoignable & Message ``Impossible de joindre l'API'' \\
\rowcolor{gray!8}
Modèle IA & Texte vide après nettoyage & ValueError explicite levée \\
\bottomrule
\end{tabular}
\end{table}

\subsection{Optimisations de performance}

\begin{itemize}
  \item \textbf{Singleton ServicePrediction :} le modèle CamemBERT (430 Mo)
        est chargé une seule fois en mémoire au démarrage. Toutes les requêtes
        partagent la même instance, évitant un rechargement à chaque appel.

  \item \textbf{torch.no\_grad() :} désactive le calcul du graphe de
        gradient lors de l'inférence, réduisant la mémoire consommée de
        50\% et améliorant la vitesse de 20\%.

  \item \textbf{Modèle en mode eval() :} \texttt{model.eval()} désactive
        le Dropout et le BatchNorm (actifs seulement à l'entraînement),
        garantissant des prédictions déterministes et plus rapides.

  \item \textbf{Pagination API :} les listes de commentaires sont paginées
        (paramètre \texttt{limite}) pour éviter de charger de grandes quantités
        de données inutilement.

  \item \textbf{IndexedStack Flutter :} évite la reconstruction des widgets
        à chaque navigation, préservant l'état et améliorant la fluidité.
\end{itemize}

\section*{Conclusion du Chapitre}

L'implémentation du système repose sur une architecture propre, modulaire
et respectueuse des bonnes pratiques. Chaque composant a été développé de
manière indépendante, communiquant via des interfaces bien définies.
Le chapitre suivant évalue les performances du système et présente les
résultats obtenus.

\cleardoublepage

% ==============================================================================
% CHAPITRE 4 : TESTS ET RÉSULTATS
% ==============================================================================
\chapter{Tests et Résultats}

\begin{infobox}[Objectif du chapitre]
Ce chapitre présente la stratégie de validation du système, les métriques
de performance du modèle d'IA, les scénarios de tests fonctionnels
et une analyse critique des résultats obtenus.
\end{infobox}

% ──────────────────────────────────────────────────────────────────────────────
\section{Stratégie de Validation}
% ──────────────────────────────────────────────────────────────────────────────

La validation du système s'articule autour de trois niveaux complémentaires :

\begin{figure}[H]
\centering
\begin{tikzpicture}[
  level/.style={rectangle, rounded corners=6pt, minimum width=10cm,
                minimum height=1cm, align=center, font=\small\bfseries,
                thick, draw},
]
\node[level, fill=bleuprincipal!15, draw=bleuprincipal] at (0,3)
  {Niveau 1 : Tests Unitaires — Fonctions individuelles (nettoyage, tokenisation, modèle)};
\node[level, fill=vertprincipal!15, draw=vertprincipal] at (0,1.5)
  {Niveau 2 : Tests d'Intégration — Communication API $\leftrightarrow$ DB $\leftrightarrow$ Modèle};
\node[level, fill=orangeaccent!10, draw=orangeaccent] at (0,0)
  {Niveau 3 : Tests Fonctionnels — Scénarios utilisateur bout en bout};
\end{tikzpicture}
\caption{Pyramide de validation du système}
\label{fig:test_pyramid}
\end{figure}

% ──────────────────────────────────────────────────────────────────────────────
\section{Performances du Modèle d'IA}
% ──────────────────────────────────────────────────────────────────────────────

\subsection{Résultats d'entraînement par époque}

\begin{table}[H]
\centering
\caption{Évolution des métriques pendant l'entraînement}
\label{tab:training_history}
\renewcommand{\arraystretch}{1.5}
\begin{tabular}{ccccc}
\toprule
\rowcolor{bleuprincipal!15}
\textbf{Époque} & \textbf{Loss Train} & \textbf{Acc. Train} &
\textbf{Loss Val.} & \textbf{Acc. Val.} \\
\midrule
1 & 0.9821 & 58.6\% & 0.8743 & 64.7\% \\
\rowcolor{gray!8}
2 & 0.6234 & 74.5\% & 0.6512 & 73.5\% \\
3 & \textbf{0.3876} & \textbf{85.4\%} & \textbf{0.5981} & \textbf{76.5\%} \\
\bottomrule
\end{tabular}
\end{table}

\subsection{Métriques finales sur le jeu de test}

Le meilleur modèle (epoch 3, loss validation minimale) est évalué sur
le jeu de test final qui n'a jamais été utilisé pendant l'entraînement :

\begin{table}[H]
\centering
\caption{Rapport de classification sur le jeu de test}
\label{tab:classification_report}
\renewcommand{\arraystretch}{1.5}
\begin{tabular}{lcccc}
\toprule
\rowcolor{bleuprincipal!15}
\textbf{Classe} & \textbf{Précision} & \textbf{Rappel} &
\textbf{F1-Score} & \textbf{Support} \\
\midrule
\textcolor{rougeprincipal}{\textbf{Négatif}} (0) & 1.0000 & 0.6364 & 0.7778 & 11 \\
\rowcolor{gray!8}
\textcolor{jauneneutre!70!black}{\textbf{Neutre}} (1) & 0.6250 & 0.8333 & 0.7143 & 12 \\
\textcolor{vertprincipal}{\textbf{Positif}} (2) & 0.8462 & 1.0000 & 0.9167 & 11 \\
\midrule
\rowcolor{bleuprincipal!8}
\textbf{Accuracy globale} & \multicolumn{4}{c}{\textbf{79.41\% (27/34)}} \\
\midrule
Macro avg & 0.8237 & 0.8232 & 0.8029 & 34 \\
\rowcolor{gray!8}
Weighted avg & 0.8253 & 0.7941 & 0.7993 & 34 \\
\bottomrule
\end{tabular}
\end{table}

\subsection{Matrice de confusion}

\begin{figure}[H]
\centering
\begin{tikzpicture}
\begin{axis}[
    axis on top,
    width=8cm, height=8cm,
    xmin=-0.5, xmax=2.5,
    ymin=-0.5, ymax=2.5,
    xtick={0,1,2},
    ytick={0,1,2},
    xticklabels={Négatif, Neutre, Positif},
    yticklabels={Positif, Neutre, Négatif},
    xlabel={\textbf{Prédit}},
    ylabel={\textbf{Réel}},
    xlabel style={font=\bfseries\small},
    ylabel style={font=\bfseries\small},
    tick label style={font=\small},
    colormap={greenblue}{
      color(0)=(white) color(1)=(bleuprincipal!80)
    },
    point meta min=0, point meta max=12
]
% Matrice : [Réel=Nég: 7,3,1], [Réel=Neu: 0,10,2], [Réel=Pos: 0,0,11]
\addplot[matrix plot*, mesh/cols=3, point meta=explicit]
coordinates {
  (0,2) [7]  (1,2) [3]  (2,2) [1]
  (0,1) [0]  (1,1) [10] (2,1) [2]
  (0,0) [0]  (1,0) [0]  (2,0) [11]
};

% Annotations
\node at (axis cs:0,2) {\large\bfseries 7};
\node at (axis cs:1,2) {\large 3};
\node at (axis cs:2,2) {\large 1};
\node at (axis cs:0,1) {\large 0};
\node at (axis cs:1,1) {\large\bfseries 10};
\node at (axis cs:2,1) {\large 2};
\node at (axis cs:0,0) {\large 0};
\node at (axis cs:1,0) {\large 0};
\node at (axis cs:2,0) {\large\bfseries 11};
\end{axis}
\end{tikzpicture}
\caption{Matrice de confusion sur le jeu de test (34 exemples)}
\label{fig:confusion_matrix}
\end{figure}

\subsection{Analyse des résultats}

\textbf{Points forts observés :}
\begin{itemize}
  \item La classe \textbf{Positif} est parfaitement reconnue (rappel = 100\%,
        F1 = 0.917). Le modèle n'a jamais manqué un commentaire positif.
  \item La classe \textbf{Négatif} atteint une précision parfaite (100\%) :
        quand le modèle dit qu'un commentaire est négatif, il a toujours raison.
  \item L'accuracy globale de \textbf{79.41\%} est remarquable compte tenu
        du jeu de données limité à 225 exemples.
\end{itemize}

\textbf{Points d'amélioration :}
\begin{itemize}
  \item La confusion principale se situe entre \textbf{Négatif et Neutre}
        (3 cas) et \textbf{Neutre et Positif} (2 cas). La frontière
        entre ces classes est naturellement floue dans le langage étudiant.
  \item Le rappel de la classe Négatif (63.64\%) est le plus faible :
        certains commentaires négatifs modérés sont classés comme neutres.
\end{itemize}

% ──────────────────────────────────────────────────────────────────────────────
\section{Tests Fonctionnels de l'API}
% ──────────────────────────────────────────────────────────────────────────────

\subsection{Scénarios de test}

\begin{table}[H]
\centering
\caption{Scénarios de test fonctionnels}
\label{tab:functional_tests}
\renewcommand{\arraystretch}{1.5}
\begin{tabular}{p{0.8cm} p{5.5cm} p{3.5cm} p{2.5cm}}
\toprule
\rowcolor{bleuprincipal!15}
\textbf{ID} & \textbf{Entrée} & \textbf{Résultat attendu} & \textbf{Statut} \\
\midrule
TF01 & ``Ce cours est excellent, j'ai tout compris'' & Positif (code 200) & \textcolor{vertprincipal}{\ding{52} Passé} \\
\rowcolor{gray!8}
TF02 & ``Cours ennuyeux, je ne comprends rien'' & Négatif (code 200) & \textcolor{vertprincipal}{\ding{52} Passé} \\
TF03 & ``Le cours se passe normalement'' & Neutre (code 200) & \textcolor{vertprincipal}{\ding{52} Passé} \\
\rowcolor{gray!8}
TF04 & ``Hi'' (2 caractères) & HTTP 422 (validation) & \textcolor{vertprincipal}{\ding{52} Passé} \\
TF05 & Texte vide \texttt{``''} & HTTP 422 & \textcolor{vertprincipal}{\ding{52} Passé} \\
\rowcolor{gray!8}
TF06 & GET /api/commentaires & Liste JSON + total & \textcolor{vertprincipal}{\ding{52} Passé} \\
TF07 & GET /api/commentaires?sentiment=2 & Uniquement Positifs & \textcolor{vertprincipal}{\ding{52} Passé} \\
\rowcolor{gray!8}
TF08 & DELETE /api/commentaires/999 & HTTP 404 & \textcolor{vertprincipal}{\ding{52} Passé} \\
TF09 & GET /api/stats (base vide) & Message ``Aucun commentaire'' & \textcolor{vertprincipal}{\ding{52} Passé} \\
\rowcolor{gray!8}
TF10 & GET /health & \texttt{\{status: ok\}} & \textcolor{vertprincipal}{\ding{52} Passé} \\
\bottomrule
\end{tabular}
\end{table}

\subsection{Exemple de réponse API}

Voici un exemple de réponse JSON retournée par \texttt{POST /api/predict}
pour le commentaire ``Ce cours est excellent, le professeur explique très bien'' :

\begin{lstlisting}[style=styleJSON, caption={Exemple de réponse JSON de l'API}, label={lst:api_response}]
{
  "id": 42,
  "texte_original": "Ce cours est excellent, le professeur explique très bien",
  "matiere": "Informatique",
  "sentiment_predit": 2,
  "label_sentiment": "Positif",
  "scores": {
    "negatif": 0.0187,
    "neutre":  0.0431,
    "positif": 0.9382
  },
  "date_creation": "2026-05-12T14:23:07"
}
\end{lstlisting}

% ──────────────────────────────────────────────────────────────────────────────
\section{Tests de l'Application Mobile}
% ──────────────────────────────────────────────────────────────────────────────

\begin{table}[H]
\centering
\caption{Scénarios de test de l'application Flutter}
\label{tab:mobile_tests}
\renewcommand{\arraystretch}{1.5}
\begin{tabular}{p{0.8cm} p{5cm} p{4cm} p{2cm}}
\toprule
\rowcolor{bleuprincipal!15}
\textbf{ID} & \textbf{Action utilisateur} & \textbf{Comportement attendu} & \textbf{Statut} \\
\midrule
TM01 & Soumettre un commentaire valide & Affichage SentimentCard colorée & \textcolor{vertprincipal}{\ding{52}} \\
\rowcolor{gray!8}
TM02 & Soumettre moins de 5 caractères & Message de validation rouge & \textcolor{vertprincipal}{\ding{52}} \\
TM03 & Analyser sans serveur démarré & Message d'erreur explicite & \textcolor{vertprincipal}{\ding{52}} \\
\rowcolor{gray!8}
TM04 & Naviguer vers Historique & Chargement de la liste & \textcolor{vertprincipal}{\ding{52}} \\
TM05 & Filtrer par Positif & Seuls les commentaires positifs & \textcolor{vertprincipal}{\ding{52}} \\
\rowcolor{gray!8}
TM06 & Glisser une carte vers la gauche & Apparition du fond rouge + corbeille & \textcolor{vertprincipal}{\ding{52}} \\
TM07 & Confirmer la suppression & Commentaire supprimé, liste rafraîchie & \textcolor{vertprincipal}{\ding{52}} \\
\rowcolor{gray!8}
TM08 & Annuler la suppression & Commentaire maintenu, liste inchangée & \textcolor{vertprincipal}{\ding{52}} \\
\bottomrule
\end{tabular}
\end{table}

% ──────────────────────────────────────────────────────────────────────────────
\section{Difficultés Rencontrées et Solutions}
% ──────────────────────────────────────────────────────────────────────────────

\begin{table}[H]
\centering
\caption{Difficultés rencontrées et solutions apportées}
\label{tab:difficulties}
\renewcommand{\arraystretch}{1.5}
\begin{tabular}{p{4cm} p{5cm} p{4cm}}
\toprule
\rowcolor{bleuprincipal!15}
\textbf{Problème} & \textbf{Cause} & \textbf{Solution} \\
\midrule
\texttt{sentencepiece} manquant & CamemBERT requiert SentencePiece pour la tokenisation & \texttt{pip install sentencepiece} \\
\rowcolor{gray!8}
Peer auth PostgreSQL & Mismatch utilisateur Linux/PostgreSQL & Modifier \texttt{pg\_hba.conf} : \texttt{peer} $\rightarrow$ \texttt{md5} \\
\texttt{cmdline-tools} Android manquant & SDK mal configuré & Téléchargement manuel + configuration \texttt{JAVA\_HOME} \\
\rowcolor{gray!8}
Flutter ne joint pas l'API & Adresse localhost incorrecte sur émulateur & Utiliser \texttt{10.0.2.2} au lieu de \texttt{127.0.0.1} \\
Gradient explosion & Taux d'apprentissage trop élevé & Gradient clipping (\texttt{max\_norm=1.0}) \\
\rowcolor{gray!8}
Déséquilibre de classes initial & Dataset non équilibré & 75 exemples par classe, stratification train/test \\
\bottomrule
\end{tabular}
\end{table}

\section*{Conclusion du Chapitre}

Les tests conduits ont validé le bon fonctionnement de l'ensemble des
composants du système. Le modèle atteint une précision de 79.41\% sur le
jeu de test, avec d'excellentes performances sur les classes Positif et
Négatif. Les quelques confusions observées entre Neutre et les classes
adjacentes sont explicables par la nature ambiguë du langage étudiant et
seront améliorées avec plus de données d'entraînement.

\cleardoublepage

% ==============================================================================
% CHAPITRE 5 : DISCUSSION ET PERSPECTIVES
% ==============================================================================
\chapter{Discussion et Perspectives}

\begin{infobox}[Objectif du chapitre]
Ce chapitre propose une analyse critique des résultats obtenus, identifie
les limites actuelles du système et trace les perspectives d'évolution
et d'amélioration futures.
\end{infobox}

% ──────────────────────────────────────────────────────────────────────────────
\section{Analyse Critique des Résultats}
% ──────────────────────────────────────────────────────────────────────────────

\subsection{Points forts du système}

\subsubsection{Performance remarquable sur un petit jeu de données}

Obtenir une précision de 79.41\% avec seulement 225 exemples d'entraînement
est un résultat significatif. Il illustre la puissance du transfer learning :
CamemBERT, ayant été pré-entraîné sur 138 Go de texte français, possède
déjà une compréhension profonde de la langue. Le fine-tuning n'a qu'à
adapter cette compréhension générale à la tâche spécifique de la
classification de sentiments étudiants.

À titre de comparaison, un modèle LSTM entraîné à partir de zéro sur le
même jeu de données obtiendrait typiquement une précision inférieure à 60\%,
car il n'aurait pas le bénéfice de ce vaste pré-entraînement.

\subsubsection{Architecture microservices modulaire}

La séparation claire entre le modèle IA, l'API, le dashboard et
l'application mobile permet :
\begin{itemize}
  \item De mettre à jour le modèle sans toucher aux interfaces utilisateur.
  \item D'ajouter de nouvelles interfaces (ex: application web) sans
        modifier le backend.
  \item De scaler chaque composant indépendamment selon la charge.
  \item De remplacer facilement une technologie par une autre (ex: MySQL
        au lieu de PostgreSQL) sans impact sur les autres couches.
\end{itemize}

\subsubsection{Accessibilité multi-canal}

Le système répond à deux profils d'utilisateurs distincts :
\begin{itemize}
  \item \textbf{L'étudiant}, qui accède via une application mobile intuitive,
        sans barrière technique.
  \item \textbf{L'enseignant}, qui dispose d'un tableau de bord analytique
        riche en visualisations.
\end{itemize}

\subsection{Limites identifiées}

\subsubsection{Taille du jeu de données}

La principale limitation du système est la taille du jeu de données
d'entraînement (225 exemples). Bien que les résultats soient encourageants,
un modèle entraîné sur 10 000 exemples ou plus serait nettement plus robuste,
notamment pour :
\begin{itemize}
  \item Les expressions idiomatiques et le langage familier étudiant.
  \item Les commentaires très courts (1-2 phrases) qui sont plus difficiles
        à classifier.
  \item Les nuances culturelles spécifiques à certains établissements.
\end{itemize}

\subsubsection{Absence d'authentification}

Le système actuel ne dispose pas d'un mécanisme d'authentification.
Tout utilisateur connaissant l'URL de l'API peut soumettre des commentaires
ou accéder aux statistiques. Dans un contexte de production réel, cela
représente une vulnérabilité importante.

\subsubsection{Déploiement local uniquement}

Le système est actuellement configuré pour fonctionner en local. Pour
un déploiement en production accessible à tous les étudiants d'un
établissement, il faudrait déployer l'API sur un serveur cloud (AWS,
GCP, Azure, ou VPS) et adapter l'URL dans l'application mobile.

\subsubsection{Pas de gestion des doublons}

Le système actuel ne détecte pas les commentaires identiques soumis
plusieurs fois, ce qui peut fausser les statistiques.

% ──────────────────────────────────────────────────────────────────────────────
\section{Impact Potentiel du Projet}
% ──────────────────────────────────────────────────────────────────────────────

Ce système, dans un contexte de déploiement réel, pourrait avoir un impact
significatif sur la qualité de l'enseignement supérieur :

\begin{table}[H]
\centering
\caption{Impact potentiel selon les parties prenantes}
\label{tab:impact}
\renewcommand{\arraystretch}{1.5}
\begin{tabular}{p{3cm} p{10cm}}
\toprule
\rowcolor{bleuprincipal!15}
\textbf{Partie prenante} & \textbf{Bénéfice concret} \\
\midrule
Étudiants & Retour d'expérience immédiat, sentiment d'être écouté, interaction plus engageante \\
\rowcolor{gray!8}
Enseignants & Identification rapide des cours mal compris, ajustement pédagogique en temps réel \\
Administration & Vue macro de la satisfaction par département, allocation optimale des ressources \\
\rowcolor{gray!8}
Établissement & Amélioration continue de la qualité, différenciation concurrentielle \\
\bottomrule
\end{tabular}
\end{table}

% ──────────────────────────────────────────────────────────────────────────────
\section{Perspectives d'Évolution}
% ──────────────────────────────────────────────────────────────────────────────

\subsection{Améliorations à court terme}

\begin{enumerate}
  \item \textbf{Agrandissement du jeu de données :} collecter 500 à 1000
        commentaires réels anonymisés d'étudiants pour re-fine-tuner le modèle
        et améliorer sa précision à 85\%+.

  \item \textbf{Authentification JWT :} implémenter JSON Web Tokens pour
        sécuriser l'API. Chaque étudiant s'authentifie et les commentaires
        sont liés à un compte anonymisé.

  \item \textbf{Interface web responsive :} ajouter une page web React ou
        Vue.js pour les étudiants ne disposant pas d'un smartphone Android.

  \item \textbf{Déploiement Docker :} conteneuriser le backend, la base de
        données et le dashboard dans des conteneurs Docker orchestrés par
        Docker Compose pour faciliter le déploiement et la portabilité.
\end{enumerate}

\subsection{Améliorations à moyen terme}

\begin{enumerate}
  \item \textbf{Analyse des aspects (Aspect-Based Sentiment Analysis) :}
        au lieu d'attribuer un seul sentiment à tout un commentaire,
        identifier les aspects mentionnés (``le contenu du cours'',
        ``la pédagogie du professeur'', ``le rythme'') et attribuer
        un sentiment à chacun.

  \item \textbf{Extraction de mots-clés :} identifier automatiquement
        les termes les plus fréquemment associés à chaque classe de sentiment
        (TF-IDF, YAKE ou KeyBERT) pour comprendre \textit{pourquoi} un
        cours est bien ou mal perçu.

  \item \textbf{Application iOS :} Flutter étant cross-platform, l'extension
        à iOS nécessite principalement une configuration Xcode et un compte
        développeur Apple.

  \item \textbf{Alertes automatiques :} envoyer un email ou une notification
        push à l'enseignant lorsque le taux de commentaires négatifs dépasse
        un seuil configurable.
\end{enumerate}

\subsection{Vision à long terme}

\begin{definitionbox}[Évolution vers un système de recommandation pédagogique]
À terme, le système pourrait évoluer vers une plateforme d'intelligence
pédagogique capable de :
\begin{itemize}
  \item Corréler les sentiments avec les résultats académiques (notes,
        taux d'abandon) pour identifier les cours à risque.
  \item Recommander automatiquement des ressources pédagogiques
        complémentaires aux étudiants dont les commentaires révèlent des
        difficultés de compréhension.
  \item Générer des rapports pédagogiques automatiques pour les conseils
        d'enseignement, alimentés par les analyses de sentiments agrégées.
  \item Intégrer des modèles multimodaux prenant en compte non seulement
        le texte mais aussi la fréquence de participation et les patterns
        temporels des commentaires.
\end{itemize}
\end{definitionbox}

\section*{Conclusion du Chapitre}

Ce chapitre a présenté une analyse honnête et nuancée du système développé.
Ses points forts résident dans l'utilisation d'une technologie de pointe
(CamemBERT), une architecture modulaire bien conçue et une accessibilité
multi-canal. Ses limites principales — liées à la taille du dataset et à
l'absence d'authentification — sont clairement identifiées et des solutions
concrètes ont été proposées. Ce travail constitue une base solide sur
laquelle des évolutions significatives pourront être construites.

\cleardoublepage

% ==============================================================================
% CONCLUSION GÉNÉRALE
% ==============================================================================
\chapter*{Conclusion Générale}
\addcontentsline{toc}{chapter}{Conclusion Générale}
\markboth{CONCLUSION GÉNÉRALE}{}

\vspace{0.3cm}

Ce mémoire a présenté la conception et le développement d'un système
complet d'analyse automatique de sentiments appliqué aux commentaires
d'étudiants sur leurs cours universitaires. Ce projet représente une
intégration réussie de plusieurs domaines de l'informatique moderne :
l'intelligence artificielle, le développement web et mobile, les bases
de données relationnelles et les architectures distribuées.

\vspace{0.5cm}

\section*{Synthèse des réalisations}

\begin{tcolorbox}[
  enhanced,
  colback=bleuprincipal!5!white,
  colframe=bleuprincipal,
  boxrule=1.5pt, arc=6pt,
  left=10pt, right=10pt, top=8pt, bottom=8pt
]
\begin{enumerate}
  \item \textbf{Modèle d'IA :} Le modèle CamemBERT fine-tuné sur 225
        commentaires étudiants en français a atteint une précision de
        \textbf{79.41\%} sur le jeu de test, démontrant la puissance du
        transfer learning pour des tâches spécialisées avec peu de données.

  \item \textbf{API REST :} Le serveur FastAPI expose 5 endpoints documentés,
        gère la validation des données via Pydantic et persiste chaque analyse
        dans PostgreSQL via SQLAlchemy ORM. Le pattern Singleton garantit
        des temps de réponse inférieurs à 1 seconde.

  \item \textbf{Dashboard analytique :} Le tableau de bord Streamlit offre
        des visualisations interactives (camembert, histogrammes, jauges)
        permettant aux enseignants de suivre en temps réel les tendances
        de satisfaction par cours et par matière.

  \item \textbf{Application mobile :} L'application Flutter offre une
        interface intuitive, matérielle et moderne permettant aux étudiants
        de soumettre leurs commentaires, consulter les résultats en temps réel
        et naviguer dans l'historique depuis n'importe quel smartphone Android.
\end{enumerate}
\end{tcolorbox}

\vspace{0.5cm}

\section*{Bilan technique}

Ce projet a permis d'explorer et de maîtriser un large spectre de
technologies : PyTorch et HuggingFace Transformers pour le deep learning,
FastAPI et SQLAlchemy pour le backend, Flutter/Dart pour le mobile, et
Streamlit/Plotly pour la visualisation. Cette polyvalence technique
constitue un acquis précieux, directement transférable dans un contexte
professionnel.

Sur le plan algorithmique, la compréhension approfondie des mécanismes
d'attention des Transformers, du fine-tuning, de l'optimiseur AdamW et
du scheduler linéaire avec warmup a enrichi notre bagage en apprentissage
automatique moderne. La résolution des défis concrets — débogage du
pipeline de données, gestion des erreurs d'authentification PostgreSQL,
configuration de l'environnement Android — a également développé
des compétences pratiques essentielles.

\vspace{0.5cm}

\section*{Perspectives immédiates}

Les trois priorités d'amélioration identifiées sont :
\begin{itemize}
  \item \textbf{Augmentation des données :} collecter 500+ commentaires
        réels pour améliorer la précision au-delà de 85\%.
  \item \textbf{Sécurisation :} implémenter l'authentification JWT et le
        déploiement Docker pour une mise en production réelle.
  \item \textbf{Analyse d'aspects :} passer d'une classification globale
        à une analyse fine des aspects pédagogiques mentionnés.
\end{itemize}

\vspace{0.5cm}

\section*{Mot final}

Ce projet de fin d'études a été l'occasion de réaliser que la frontière
entre recherche et application concrète est plus mince qu'il n'y paraît.
CamemBERT, un modèle sorti de laboratoires de recherche de renom, peut
être adapté en quelques jours pour résoudre un problème réel et local.
Cette capacité à connecter la recherche fondamentale à des besoins
pratiques constitue, à notre sens, l'essence même de l'ingénierie
informatique moderne.

\vspace{0.5cm}

\begin{flushright}
\textit{``La technologie est à son meilleur quand elle rapproche les gens.''}\\
\vspace{0.3cm}
\textbf{Hissein Nasser} — Mai 2026
\end{flushright}

\cleardoublepage


% ── Bibliographie ─────────────────────────────────────────────────────────────
\printbibliography[heading=bibintoc, title={Références Bibliographiques}]

\end{document}
